\documentclass[
%%% one of
%submission
final
%proceedings
%%% if you compile a final version for the old OJS platform
% , ojs
%%% if all authors have the same affiliation
% , nomarks
]{dmtcs-episciences}

% DON'T LOAD ANY STYLES THAT CHANGE THE PAGE LAYOUT
% AND DON'T CHANGE THE PAGE LAYOUT BY HAND, EITHER.


\usepackage[utf8]{inputenc}
\usepackage{subfigure}

% graphicx is now loaded automatically no need to put this in here anymore.
%
%\usepackage{graphicx}

% We strongly recommend to use natbib. Your colleagues deserve to be
% named in your text. PLEASE, ADAPT YOUR TEXT ACCORDINGLY, such that
% citations are grammatically correct.
%\usepackage[round]{natbib}

\author{Alexander Burstein\affiliationmark{1}
  \and Opel Jones\affiliationmark{2} 
  }

\title{Enumeration of Dumont permutations avoiding certain four-letter patterns}
% put your affiliation here, not your full address.
% If you like to give away your email or other parts of your address,
% THIS IS NOT THE RIGHT PLACE, your address will change, this paper
% will not.
% Just watch that your personal data that you want to communicate on
% the episcience server is always up to date.

\affiliation{
  % one line per affiliation, no postal codes, grant numbers or similar
  Department of Mathematics, Howard University, Washington, DC 20059, USA\\
  Department of Mathematics, Towson University, Towson, MD 21252, USA
  }

\keywords{Dumont permutation, permutation pattern, pattern avoidance}
% don't try to cheat here, we will check the dates!

\received{2020-02-28}
\revised{2021-05-07}
\accepted{2021-05-12}
\publicationdetails{22}{2021}{2}{7}{6174}

%\maketitle


%\documentclass[11pt]{article}

\usepackage{amsthm,amsfonts,amssymb,amsopn,
%fullpage,
amsmath,indentfirst,tikz,hyperref,varioref}
%\usepackage[margin=1.2in]{geometry}
\usepackage{extpfeil}

\newtheorem{theorem}{Theorem}[section]%[section]
%\newtheorem{theorem2}{Theorem2} %[for different numbering style]
\newtheorem{prop}[theorem]{Proposition}
\newtheorem{lemma}[theorem]{Lemma}
\newtheorem{corollary}[theorem]{Corollary}
\newtheorem{conjecture}[theorem]{Conjecture} 

{
\theoremstyle{definition}
\newtheorem{definition}[theorem]{Definition}
\newtheorem{notation}[theorem]{Notation}
\newtheorem{example}[theorem]{Example}
}

{
\theoremstyle{remark}
\newtheorem{remark}[theorem]{Remark}
%\newtheorem{example}[theorem]{Example}
}

\DeclareMathOperator{\Av}{Av}
\DeclareMathOperator{\pk}{pk}
\DeclareMathOperator{\pt}{pt}
%\setcounter{section}{-1}

%\renewcommand{\baselinestretch}{2}

%\author{
%\textsc{Alexander Burstein}\\
%Department of Mathematics\\
%Howard University\\
%Washington, DC 20059\\
%\texttt{aburstein@howard.edu}
%\and
%\textsc{Opel Jones}\\
%Department of Mathematics\\
%Towson University\\
%Towson, MD 21252\\
%\texttt{ojones@towson.edu}
%}
%
%\title{\textbf{Enumeration of Dumont permutations avoiding certain four-letter patterns}}
%
%\date{May 5, 2021}
%
%\maketitle


\begin{document}

\maketitle

%%%%%%%%%%%%%%%%%%%%
% Begin Abstract

\begin{abstract}
In this paper, we enumerate Dumont permutations of the fourth kind avoiding or containing once certain permutations of length 4. We also conjecture a Wilf-equivalence of two 4-letter patterns on Dumont permutations of the first kind.
\end{abstract}

%
%\begin{abstract}
%
%Pattern avoidance in permutations and minimal occurrences (one or two) of a pattern in permutations are fairly new subjects; each are less than fifty years young.  In 1974 Dumont found two types of permutations are counted the same sequence. The first type is a permutation in which each even entry is followed by a smaller entry, and each odd entry is followed by a larger entry, or ends the permutation. The second type is a permutation wherein if an entry is a deficiency, it must be even, and if an entry is an excedance or a fixed point, it must be odd. These are now known as Dumont permutations of the first and second kinds. 
%
%Since then, Burstein and Stromquist \cite{Burstein_Third} and Burstein, Josuat-Verg\`es, and Stromquist \cite{Burstein_New} found two other types of permutations which are also counted by the same sequence. The third type is a permutation in which each descent is from an even entry to an even entry.  The fourth type is a permutation in which deficiencies must be even entries, and can only occur at even positions.  These are now known as Dumont permutations of the third and fourth kinds
%
%In this paper we will discuss several enumerations of restricted Dumont permutations, that is Dumont permutations avoiding certain patterns. We will also briefly discuss their proofs which involve methods using induction, block decomposition, Dyck paths, and generating functions.  We will conclude with a conjecture that the patterns 2143 and 3421 are indeed Wilf-equivalent on Dumont permutations of the first kind.  We were fortunate enough to receive the assistance of Michael Albert with additional computations which strengthened our conjecture for $n$ up to 10, that is Dumont permutations of the first kind of length up to 20.
%
%\end{abstract}

% End Abstract

%%%%%%%%%%%%%%%%%%%%

% Begin Section 1 - Preliminaries

\section{Preliminaries}

This paper is concerned with the enumeration of pattern-restricted Dumont permutations. We will begin by defining patterns and Dumont permutations of various kinds, the original two kinds defined by Dumont \cite{Dumont} and the more recent third and fourth kinds defined by Burstein et al. \cite{Burstein_New}. We will then discuss previous results on this topic and present several new avoidance results for Dumont permutations of the fourth kind, an enumeration result for Dumont permutations of the fourth kind containing a single occurrence of a pattern, as well as a Wilf-equivalence conjecture on restricted Dumont permutations of the first kind.

\subsection {Patterns}

%The notion of pattern avoidance is fairly new and a rapidly growing subject in the realm of combinatorics.  Many interesting results connect permutation patterns to various parts of mathematics including enumerative and algebraic combinatorics, probability theory, representation theory, and algebraic geometry. At the same time, permutations patterns are used outside of mathematics in subjects such as theoretical computer science and computational biology.  Multiple open questions are often brought up in various conferences and papers on the subject.

We begin with an example of pattern containment.  Suppose we are given a permutation, say $p=26483751$, and another permutation usually of shorter or equal length, e.g. $q=132$.  We say that the string $265$ is \emph{order-isomorphic} to $132$.  In such a case, we say that $p$ contains $q$ as a \emph{pattern}, and that $265$ is an \emph{instance}, or \emph{occurrence}, of $q$ in $p$.  Notice also that $265$ is not the only instance of $q$ in $p$, so are the subsequences $287$ and $475$. 

On the other hand, for the same $p$ and a pattern $q'=1234$,  we see that there is no length $4$ subsequence of $p$ that is order-isomorphic to $q'$. In this case we say that $p$ \emph{avoids} $q'$.  Let us now define pattern containment and pattern avoidance more formally.

\subsection{Pattern avoidance}

Simion and Schmidt \cite{Simion} was the first paper devoted solely to pattern avoidance. Since then, pattern avoidance has been the subject of numerous papers. We will only mention here a few introductory sources:
\begin{itemize}

\item Bevan \cite{Bevan} contains essential terminology used in many patterns on permutations patterns,

\item  Chapter 4 of B\'ona's \emph{Combinatorics of Permutations} \cite{Bona} contains a textbook introduction to the subject, 

\item Vatter's ``Permutation Classes'' \cite{Vatter} in \emph{Handbook of Enumerative Combinatorics} is an in-depth survey of the literature and foundational results. 

\end{itemize}

%, including e.g. \cite{Albert, Bona_length4, Burstein_Restricted, Burstein_Dyck, Burstein_New, Burstein_Third, Mansour, MV_Restricted, MV_Chebyshev, Noonan_Zeilberger, Stankova}.

%\begin{notation}
%Let $S_n$ be the set of permutations of $[n] = \{1,2,3,...,n\}$.
%\end{notation}

\begin{definition}[\cite{Bona}]
\label{def:avoidance}
Let $S_n$ be the set of permutations of $[n] = \{1,2,3,...,n\}$. Let $p=(p_1,p_2,...,p_n)\in S_n$ and $q=(q_1,q_2,...,q_k)\in S_k$ be two permutations.  Then we say that $p$ \emph{contains} $q$ if there is a $k$-tuple $(p_{i_1},p_{i_2},\dots,p_{i_k})$ in $p$ such that $1\le i_1<i_2<\dots<i_k\le n$, and $p_{i_\alpha}<p_{i_\beta}$ if and only if $q_\alpha<q_\beta$.  Otherwise, we say that $p$ \emph{avoids} $q$.  
\end{definition}

\begin{notation}
For a set of permutations avoiding a pattern or set of patterns, we let $S_n(\tau)$ be the set of permutations in $S_n$ avoiding a pattern $\tau$, and we let $S_n(T)$ be the set of permutations avoiding a set of patterns $T$, in other words, simultaneously avoiding all patterns in $T$.
\end{notation}

%Note that in Definition \ref{def:avoidance} we will usually have $k \le n$. We also need the following notation in showing pattern avoidance for theorems in Subsection \ref{length four_thms}.

\begin{notation}
Let $\tau'$ and $\tau''$ be subsequences of a permutation $\tau$.  We say that $\tau' > \tau''$ if every entry of $\tau'$ is greater than every entry of $\tau''$.  We define $\tau' < \tau''$ similarly.
\end{notation}

\begin{notation}
For a string $\tau$ and an integer $m$, the string $\tau + m$ (resp. $\tau - m$) is obtained by adding $m$ to (resp. subtracting $m$ from) every entry of $\tau$.
\end{notation}

Thus, for example, if $\tau$ is a permutation of $[n]$, then $\tau+m$ is a permutation of $[m+1,m+n]$.

Since $|S_n(12)| = |S_n(21)| = 1$ for all $n \ge 0$, we are interested in patterns of length $n \ge 3$. For permutations avoiding patterns of length $n=3$, Knuth \cite{Knuth} and Simion and Schmidt \cite {Simion} showed that permutations avoiding any $\tau \in S_3$ are Wilf-equivalent, i.e. have the same enumeration sequence.  More precisely, for all $\tau\in S_3$ and $n \ge 0$, $|S_n(\tau)| = C_n$, where $C_n$ is the $n$-th Catalan number. 

%We will expound upon the aforementioned enumeration after exploring permutation symmetries as well as Wilf-equivalence.

\subsection{Permutation symmetries and Wilf-equivalence}

%To assist with showing that all permutations avoiding patterns of length three are enumerated by the Catalan numbers, we will use 

Permutation symmetries are quite convenient, in that they help us find equinumerously avoided patterns (or sets of patterns, respectively). 

\begin{definition}
Let patterns $\tau_1$ and $\tau_2$ be such that $|S_n(\tau_1)| = |S_n(\tau_2)|$ for any $n \ge 0$.  Also, let sets of patterns $T_1$ and $T_2$ be such that $|S_n(T_1)| = |S_n(T_2)|$ for any $n \ge 0$.  Then such patterns (or sets of patterns, respectively) are called \emph{Wilf-equivalent} and said to belong to the same \emph{Wilf class}.
\end{definition}

\begin{definition}
Let $\pi \in S_{n}$ be a permutation, and let $\pi(i)$ be the $i$-th entry of $\pi$ from left to right.  Define a \emph{permutation diagram of $\pi$} as the set of dots $\{(i, \pi(i))\, |\, 1 \le i \le n\}$.
\end{definition}

Visually, we can represent $\pi$ by placing dots on an $n \times n$ board in the following fashion: place a dot inside the square in the $i$-th column from the left and $\pi(i)$-th row from the bottom.  Thus, we read $\pi$ from left-to-right and bottom-to-top, as mentioned earlier.  Note, the origin of our $n \times n$ board is at the bottom-left corner.  %See Figure \ref{fig:big8x8_1342} on page \pageref{fig:big8x8_1342}.

%See Figure \ref{fig:big6x6}.
%
%\begin{example}
%
%Let $\pi = 263541 \in S_6$.  The corresponding permutation diagram of $\pi$ on a $6 \times 6$ board is
%
%\begin{center}
%\begin{tikzpicture} 
%\draw[scale = 0.5, step = 2, very thin, color=gray] (0,0) grid (12,12);
% \foreach \position in {(1,3), (3,11), (5,5), (7,9), (9,7), (11,1)}
%    \filldraw [scale = 0.5, color = black] \position circle (8pt);
%\end{tikzpicture} 
%\begin{figure}[ht]
%\caption{Permutation diagram of $\pi = 263541$  on a $6 \times 6$ board}
%\label{fig:big6x6}
%\end{figure}
%\end{center} 
%
%\end{example}
%
%%When convenient, instead of displaying a permutation on an $n \times n$ board, we will use an $n \times n$ grid wherein dots are placed on the intersections, similar to the first quadrant of a Cartesian coordinate system, in the following fashion: the abscissa relates to $i$ and the ordinate to $\pi(i)$.  Again, the origin of the grid is at the bottom-left corner.  See Figure \ref{fig:big6gx6g}.
%
%%\newpage
%
%\begin{center}
%\begin{tikzpicture} 
%\draw[scale = 0.5, step = 2, thin, color=gray] (-1,-1) grid (11,11);
% \foreach \position in {(0,2), (2,10), (4,4), (6,8), (8,6), (10,0)}
%    \filldraw [scale = 0.5, color = black] \position circle (8pt);
%\end{tikzpicture} 
%\begin{figure}[ht]
%\caption{Permutation diagram of $\pi = 263541$ on a $6 \times 6$ grid}
%\label{fig:big6gx6g}
%\end{figure}
%\end{center} 

Now let us define a few symmetry operations on $S_n$ that map every pattern onto a Wilf-equivalent pattern on the ambient set $S_n$.

\begin{definition}
Let $\pi \in S_n$ be a permutation, read left-to-right and bottom-to-top.

\begin{itemize}
\item The \emph{reverse} of $\pi$, denoted $\pi^\text{\hspace{0.01in}r}$, is $\pi$ read right-to-left and bottom-to-top.
\item The \emph{complement} of $\pi$, denoted $\pi^\text{\hspace{0.01in}c}$, is $\pi$ read left-to-right but top-to-bottom.
\item The \emph{composition of the reverse and the complement} of $\pi$, denoted  $\pi^\text{\hspace{0.01in}r\hspace{0.01in}c}$, is $\pi$ read right-to-left and top-to-bottom. Note that $\pi^\text{\hspace{0.01in}r\hspace{0.01in}c} = \pi^\text{\hspace{0.01in}c\hspace{0.01in}r}$.
\item The \emph{inverse} of $\pi$, denoted $\pi^{-1}$, is the usual inverse of a permutation.  That is, the position of each entry and the entry itself of $\pi$ are switched in $\pi^{-1}$.
\end{itemize}

\end{definition}

%
%\begin{notation}
%With respect to each specific entry,
%\begin{itemize}
%\item reverse: $\pi^\text{\hspace{0.01in}r} (i) = \pi (n + 1 - i)$.
%\item complement: $\pi^\text{\hspace{0.01in}c} (i) = n + 1 - \pi ( i)$.
%\item composition of reverse and complement: $\pi^\text{\hspace{0.01in}r\hspace{0.01in}c} (i) = \pi^\text{\hspace{0.01in}c\hspace{0.01in}r} (i) = n + 1 - \pi (n + 1 - i)$.
%\item inverse: $\pi^{-1}(i) = j$ if and only if $\pi(j) = i$.
%\end{itemize}
%\end{notation}


\begin{definition}

The symmetries above generate the group of symmetries of the square and yield the set of patterns
$\{ 
\pi, 
\pi^\text{\hspace{0.01in}r}, 
\pi^\text{\hspace{0.01in}c}, 
\pi^\text{\hspace{0.01in}r\hspace{0.01in}c}, 
\pi^{-1},
(\pi^{-1})^\text{\hspace{0.01in}r}, 
(\pi^{-1})^\text{\hspace{0.01in}c}, 
(\pi^{-1})^\text{\hspace{0.01in}r\hspace{0.01in}c}
\}$,
called the \emph{symmetry class} of $\pi$.

\end{definition}


\begin{example}

%\begin{itemize}
%\item $\pi^\text{\hspace{0.01in}r} = 145362$ 
%\item $\pi^\text{\hspace{0.01in}c} = 514236$
%\item $\pi^\text{\hspace{0.01in}r\hspace{0.01in}c} = 632415$
%\item $\pi^{-1} = 613542$
%\item $(\pi^{-1})^\text{\hspace{0.01in}r} = 245316$
%\item $(\pi^{-1})^\text{\hspace{0.01in}c} = 164235$ 
%\item $(\pi^{-1})^\text{\hspace{0.01in}r\hspace{0.01in}c} = 532461$
%\end{itemize}
%
%\end{example}

Let $\pi = 263541$.  Then the symmetry class of $\pi$ is the set of patterns
\[
\{263541,\ 145362,\ 514236,\ 632415,\ 613542,\ 245316,\ 164235,\ 532461\},
\]
\end{example}

Note the following property of the symmetry classes of permutations: if $\pi$ avoids a pattern $\rho$, then each symmetry of $\pi$ avoids the same symmetry of $\rho$.  In other words, if $\pi$ avoids $\rho$, then $\pi^\text{\hspace{0.01in}r}$ avoids $\rho^\text{\hspace{0.01in}r}$,
$\pi^\text{\hspace{0.01in}c}$ avoids $\rho^\text{\hspace{0.01in}c}$, and so on.
%$\pi^\text{\hspace{0.01in}r\hspace{0.01in}c}$ avoids $\rho^\text{\hspace{0.01in}r\hspace{0.01in}c}$,
%$\pi^{-1}$ avoids $\rho^{-1}$, and so on.
Thus, each element of the symmetry class of $\pi$ is Wilf-equivalent to $\pi$.

%
%
%\begin{example}
% Continuing with our running example, let $\pi$ yet again be $263541$, and let $\rho = 213$.  We see that there is no 3-tuple in $\pi$ that is order-isomorphic to $\rho$.  That is, $\pi$ avoids $\rho$.  Then,
%\begin{itemize}
%\item $\pi^\text{\hspace{0.01in}r} = 145362$ avoids $\rho^\text{\hspace{0.01in}r} = 312$,
%\item $\pi^\text{\hspace{0.01in}c} = 514236$ avoids $\rho^\text{\hspace{0.01in}c} = 231$,
%\item $\pi^\text{\hspace{0.01in}r\hspace{0.01in}c} = 632415$ avoids $\rho^\text{\hspace{0.01in}r\hspace{0.01in}c} = 132$,
%\item $\pi^{-1} = 613542$ avoids $\rho^{-1} = 213$, and so on.
%\end{itemize}
%\end{example}
%
%
%\subsection{Permutation statistics} 
%
%Let us now discuss a few permutation statistics.  They are indeed necessary for this paper in terms of the definitions of Dumont permutations of the first, second, third, and fourth kinds.
%
%\begin{definition}
%\cite {Burstein_Statistics}
%A combinatorial \emph{ statistic} on a set $S$ is a map $\textbf{f}: S \rightarrow \mathbb{N}^m$ for some integer $m \ge 0$.  
%\end{definition}
%
%\begin{definition} 
%\cite {Burstein_Statistics}
%The \emph{distribution} of \textbf{f} is the map $\text{d}_{\textbf{f}} :  \mathbb{N}^m \rightarrow \mathbb{N}$ with $\text{d}_{\textbf{f}}(\textbf{i}) = |\textbf{f}^{\hspace{0.05in}-1}(\textbf{i})|$ for $\textbf{i} \in \mathbb{N}^m$, where $ |\textbf{f}^{\hspace{0.05in}-1}(\textbf{i})|$ is the number of objects $s \in S$ such that $\textbf{f}(s) = \textbf{i}$. 
%\end{definition}
%
%\begin{definition}
%\cite {Bursten_Mahonian}
%We say that statistics \textbf{f} and \textbf{g} are \emph{equidistributed} and write $\textbf{f} \sim \textbf{g}$ if $\text{d}_{\textbf{f}} = \text{d}_{\textbf{g}}.$
%\end{definition}
%
%There are several well known combinatorial statistics on $S_n$ that we will need for on to define Dumont permutations. For Dumont permutations of the first and third kinds, we define descents and ascents. For Dumont permutations of the second and fourth kinds, we define fixed points, excedances, and deficiencies.
%
%\begin{definition}
%We say that $i\in[n-1]$ is a \emph{descent} of $\pi$ if $\pi(i) > \pi(i+1)$.  
%In this case, we call $\pi(i)$ the \emph{descent top} and $\pi(i+1)$ the \emph{descent bottom}. We also say that $i\in[n-1]$ is an \emph{ascent} of $\pi$ if $\pi(i) < \pi(i+1)$. In this case, we call $\pi(i)$ the \emph{ascent bottom} and $\pi(i+1)$ the \emph{ascent top}.
%\end{definition}
%%
%\begin{example}
%\label {descent}
%Let $\pi = 243165 \in S_6$. Then positions 2, 3, and 5 are descents of $\pi$, while positions 1 and 4 are ascents of $\pi$.  Note that the last position, in this case occupied by the entry 5, is neither a descent nor an ascent.
%\end{example}
%
%For Dumont permutations of the second and fourth kinds, we define fixed points, excedances, and deficiencies.

\begin{definition}
A \emph{fixed point} (resp. \emph{excedance}, \emph{deficiency}) of a permutation $\pi$ is a position $i$ such that $\pi(i) =  i$ (resp. $\pi(i) > i$, $\pi(i) <  i$).
\end{definition}
%
%\begin{example}
%Let $\rho = 28145736 \in S_8$. Then:
%\begin{itemize}
%\item the fixed points of $\rho$ are 4, and 5,
%\item the excedances of $\rho$ are 2, 8, and 7,
%\item the deficiencies of $\rho$ are 1, 3, and 6.
%\end{itemize}
%\end{example}
%%
%\begin{definition}
%\cite{Bona}
%Let $\pi = \pi_1 \pi_2 \cdots \pi_{n-1} \pi_n$ be a permutation.
%An \emph{inversion} of $\pi$ is a pair $(\pi_i,\pi_j)$ whenever $i < j$ but $\pi_i > \pi_j.$
%\end{definition}
%
%\begin{example}
%Let $\pi = 243165$ as in Example \ref {descent}. Then the five inversions of $\pi$ are $(2, 1), (4, 3), (4, 1), (3, 1),$ and $(6, 5).$
%\end{example}

%
%\subsection{Permutations avoiding patterns of length three}
%
%Now that we have defined a few permutation statistics needed for Dumont permutations, as well as a few permutation symmetries needed for the enumeration of permutations avoiding patterns of length three, we are prepared to show, via Simion and Schmidt \cite {Simion}, that all permutations avoiding $\tau \in S_3$ enumerate the same sequence; the Catalan numbers.  More specifically, for all $n \ge 0$, $|S_n(\tau)| = C_n$, where $C_n$ is the $n$-th Catalan number.
%
%To begin, there are only $3! = 6$ patterns of length three in $S_3$, namely 123, 132, 213, 231, 312, and 321.  Given our earlier discussion, we know that for any permutation avoiding 123, its reverse would avoid 321. Thus, $|S_n(123)| = |S_n(321)|$.  Note that 123 and 321 have other symmetries in addition to reversal.  That is 123 and 321 are complements of each other, and 123 and 321 are their own respective inverses.
%
%Now consider 132.  Again, for any permutation avoiding 132, its reverse avoids 231, its complement avoids 312, and its reverse complement (or complement reverse) avoids 213.  Thus, $|S_n(132)| = |S_n(231)| = |S_n(312)| = |S_n(213)|$.  If it can be shown that $|S_n(123)| = |S_n(132)|$, we will be one step closer to our desired result.  Let us now define a left-to-right minimum.
%
%\begin{definition}
%\label {LR Minimum}
%A \emph{left-to-right minimum} is any entry of a permutation that is smaller than all entries which precede it, reading left to right.  Note that the first entry of a permutation is always a left-to-right minimum.
%\end{definition}
%
%We will now reconstruct the proof by Simion and Schmidt which shows that $S_n(123)$ and $S_n(132)$ are equinumerous.
%
%\begin{lemma}
%\cite {Simion}
%For all positive integers n, $|S_n(123)| = |S_n(132)|$.
%\end{lemma}
%
%\begin{proof}
%Let us construct a bijection $f$ from the set of 132-avoiding permutations to the set of 123-avoiding permutations.  Let $\pi$ be a permutation that is 132-avoiding. The map $f$ is as follows:
%
%\begin{itemize}
%\item Keep the left-to-right minima of $\pi$ fixed in place.
%\item Take all of the remaining entries which are not left-to-right minima, and reorder them from left to right in decreasing order. 
%\item This new permutation, $f(\pi)$, is 123-avoiding.  Note, $f(\pi)$ is the union of two decreasing subsequences, one of which consists of left-to-right minima, and the other consists of the remaining entries. 
%\end{itemize}
%
%\begin{example}
%Let $\pi = 57361248$.  The left-to-right minima of $\pi$ are 5, 3, and 1; leave them fixed in place.  Take the remaining entries and reorder them from left to right in decreasing order.  Thus $f(\pi) = 58371642$.
%\end{example}
%
%Comparing $\pi$ and $f(\pi)$, one can see that the left-to-right minima are the same.  In other words, rearranging the entries from $\pi$ that are not left-to-right minima did not change the left-to-right minima in $f(\pi)$.  Also, note that $f(\pi)$ is a unique 123-avoiding permutation with the given entries and positions of left-to-right minima.  Indeed, if there were two entries $\rho_j$ and $\rho_k$ that are not left-to-right minima and form a 12-subsequence, then the left-to-right minimum $\rho_i$ on the left of $\rho_j$, would make an increasing subsequence, specifically a 123-subsequence. Thus, $f(\pi)$ would not be 123-avoiding.
%
%Now, to show the inverse of the map $f$, let $\rho$ be a permutation that is 123-avoiding.  Let's call this map $g$, and it is as follows:
%
%\begin{itemize}
%\item Keep the left-to-right minima of $\rho$ fixed in place.
%\item Take all of the remaining entries which are not left-to-right minima, and reorder them from left to right in the following fashion:
%
%\begin{itemize}
%\item Fill the leftmost entry with the smallest element which is larger than the nearest left-to-right minimum on the left of that position.  
%\end{itemize}
%
%\item This new permutation, $g(\rho)$, is 132-avoiding.
%\end{itemize}
% 
%\begin{example}
%Let $\rho = 58371642$.  The left-to-right minima of $\pi$ are 5, 3, and 1; leave them fixed in place.  Take the remaining entries and reorder them from left to right like so:
%
%\begin{itemize}
%\item In the empty position between 5 and 3, place the smallest entry closest to 3 which is larger than 5, that is 7.
%\item In the empty position between 3 and 1, place the smallest entry closest to 1 which is larger than 3, that is 6.
%\item In the empty position on the right of 1, place the smallest entry not yet used, which is larger than 1, that is 2.
%\item Complete the algorithm by placing 4 and 8 in the remaining positions.
%\end{itemize}
%
%Thus $g(\rho) = 57361248$.
%\end{example}
%
%We see that $g(\rho)$ is 132-avoiding permutation, and again, that is a unique 132-avoiding permutation with the given entries and positions of left-to-right minima.  Indeed, if $g(\rho)$ contained an instance of 132, then that instance would begin with a left-to-right minimum, which cannot happen since entries larger than any left-to-right minima form an increasing subsequence.  Also, if at any given instance, two entries $\pi_j < \pi_k$ that are larger than a left-to-right minimum $\pi_i$, then $\pi_i \pi_j \pi_k$ would form a 132 subsequence. 
%
%Finally, this proves the lemma, wherein $g(f(\pi)) = \pi$ implies that $f$ is a bijection.
%\end{proof}
%
%Now to prove that all permutations avoiding any pattern of length three enumerate the Catalan numbers, that is for all $n \ge 0$, $\tau \in S_3$, $|S_n(\tau)| = C_n$, where $C_n$ is the $n$-th Catalan number.  It is safe to choose any pattern of length three.  Again, we reconstruct the proof by Simion and Schmidt using the pattern 132.
%
%\begin{theorem}
%\cite {Simion}
%For all $n \ge 0$, we have 
%$$|S_n(132)| = C_n = \frac{1}{n+1} \binom{2n}{n}.$$
%\end{theorem}
%
%\begin{proof}
%To start, let $c_n = S_n(132)$.  Suppose we have a 132-avoiding permutation wherein the entry $n$ is in the $i$-th position.  Since it avoids the pattern 132, any entry to the left of $n$ will be larger than any entry to the right of $n$.  Thus, all entries to the left of $n$ must belong to the set $\{n-i+1, n-i+2, ..., n-2,n-1\}$, and all entries to the right of $n$ must belong to the set $[n-i]$.  Now these strings, which themselves avoid 132, must be $c_{i-1}$ and $c_{n-i}$ in number, to the left of $n$ and to the right of $n$, respectively.  Summing for all $i$, the following recursion arises:
%$$c_n = \sum_{i=0}^{n-1} c_{i-1} \, c_{n-i}\, .$$
%
%Now, considering the ordinary generating function, $C(x) = \sum_{n=0}^{\infty} c_n \,  x^n$, using a convenient method from \cite{Wilf}, our recursion implies that  $$C(x)=1+xC^2(x) \, ,$$ which yields
%$$C(x) = \frac{1 - \sqrt{1 - 4x}}{2x} \, ,$$
%
%\noindent which further yields 
%$$C(x) = \sum_{n=1}^{\infty} \frac {1}{n} \, \binom{2n-2}{n-1}  \,  x^{n-1}  = \sum_{n=0}^{\infty} \frac {1}{n+1} \, \binom{2n}{n}  \,  x^{n} \, ,$$
%
%\noindent and the theorem is proved.
%\end{proof}
%
%Lastly, since $S_n(132)$ is equinumerous to $S_n(\upsilon)$, for any $\upsilon$ of length three, we have the following corollary:
%
%\begin{corollary}
%\cite{Simion}
%Let $\tau$ be a pattern of length three. Then,
%$$|S_n(\tau)| = C_n = \frac{1}{n+1} \binom{2n}{n} \, .$$
%\end{corollary}
%
%\subsection{Permutations avoiding patterns of length four}
%
%Let us now take a look at permutations avoiding patterns of length four.  To begin, there are $4! = 24$ patterns of length four in $S_4$, namely:
%\begin{center}
%$1234 \qquad 1243 \qquad 1324 \qquad 1342 \qquad 1423 \qquad 1432$
%$2134 \qquad 2143 \qquad 2314 \qquad 2341 \qquad 2413 \qquad 2431$
%$3124 \qquad 3142 \qquad 3214 \qquad 3241 \qquad 3412 \qquad 3421$
%$4123 \qquad 4132 \qquad 4213 \qquad 4231 \qquad 4312 \qquad 4321$
%\end{center}
%
%Taking reverses, complements, and inverses, and then any compositions, yields seven symmetry classes, namely:
%\begin{center}
%$1234 \qquad 1243 \qquad 1324 \qquad 1342 \qquad 1432 \qquad 2143 \qquad 2413.$
%\end{center}
%
%Now to cut down the list a little further, we use a theorem by Backelin, West, and Xin.
%
%\begin{theorem}
%\label{Backelin}
%\cite{Backelin} Let $k$ be any positive integer, and let $p$ be a permutation of the set $\{k+1, k+2, \, \cdots, k+r\}$.  Then for all positive integers n, we have
%$$S_n (123 \cdots kp) = S_n (k(k-1) \cdots 1p).$$
%\end{theorem}
%
%\begin{example}
%Using Theorem \ref{Backelin}, we have the following:
%\begin{itemize}
%\item If $r = 2$, $k = 2$, and $p = 34$, then $|S_n (1234)| = |S_n (2134)|$.  Since 2134 is the reverse-complement of 1243, we can extend to $|S_n (1234)| = |S_n (2134)| = |S_n (1243)|$.
%\item If $r = 2$, $k = 2$, and $p = 43$, then $|S_n (1243)| = |S_n (2143)|$.
%\item If $r = 3$, $k = 1$, thus $p = 4$, then $|S_n (1234)| = |S_n (3214)|$.  Again, since 3214 is the reverse-complement of 1432, we can extend to $|S_n (1234)| = |S_n (3214)| = |S_n (1432)|$.
%\end{itemize}
%\end{example}
%
%In other words, the number of permutations avoiding $1234$ is equal to the number of permutations avoiding $1243, 2143$, and $1432$.  We now have only four patterns, namely:
%\begin{center}
%$1234 \qquad 1324 \qquad 1342 \qquad 2413.$
%\end{center}
%
%The work of Stankova will get rid of one more pattern.
%
%\begin{theorem}
%\label{Stankova}
%\cite{Stankova}
%For all positive integers $n$, we have
%$$|S_n (4132)| = |S_n (3142)|.$$
%\end{theorem}
%
%As the pattern $4132$ is the complement-inverse of $1342$ and the pattern $3142$ is the reverse of $2413$, via Theorem \ref{Stankova} we have $|S_n (1342)| = |S_n (2413)|$ which leaves us with only three patterns, namely:
%\begin{center}
%$1234 \qquad 1324 \qquad 1342.$
%\end{center}
%
%To date, these three patterns result in the distinct cardinality classes of permutations avoiding patterns of length four, namely $S_n (1234), S_n (1324),$ and $S_n (1342)$.  Their respective enumerations are not trivial.  The enumeration of $|S_n (1234)|$ is an exact formula:
%
%\begin{theorem}
%\cite {Gessel}
%$$|S_n (1234)| = 2 \cdot \sum_{k=0}^n \, \binom{2k}{k} \, \binom{n}{k}^2 \frac{3k^2 + 2k + 1 - n - 2nk}{(k+1)^2 (k+2) (n - k + 1)}.$$
%\end{theorem}
%
%\noindent We use the last result to compare with $|S_n (1324)|$: 
%
%\begin{theorem}
%\cite{Bona_length4} 
%For all $n > 7$, we have, 
%$$|S_n(1234)| < |S_n(1324)|,$$
%\end{theorem}
%
%\noindent which B\'ona conjectures has an an upper bound, most likely to be decreased:
%
%\begin{conjecture}
%\cite{Bona} 
%For all positive integers $n$, we have, 
%$$|S_n(1324)| < 9^n.$$
%\end{conjecture}
%
%\noindent Lastly, the enumeration of $|S_n (1342)|$ also has an exact formula:
%
%\begin{theorem}
%\cite{Bona} 
%For all positive integers $n$, we have, 
%$$|S_n(1342)|  = (- 1)^{n - 1} \cdot \frac{(7n^2 - 3n - 2)}{2} + 3 \sum_{i = 2}^n \, (- 1)^{n - i} \cdot 2^{i+1} \cdot \frac{(2i - 4)!}{i!(i - 2)!} \cdot \binom{n - i + 2}{2}, $$
%\end{theorem}
%
%\noindent which has an upper bound, namely $|S_n(1342)| < 8^n.$
%
%
%\subsection{Permutations with one or two pattern occurrences}
%
%We now take a look at permutations with one or two pattern occurrences.  The earliest papers to analyze one or two occurrences of a pattern were less than a quarter century ago; this distinction goes to Noonan, and Noonan and Zeilberger \cite {Noonan, Noonan_Zeilberger}, as they were the first to truly review this phenomenon.  A short time thereafter, Bona \cite {Bona_132} expounded upon their work.  Since then, the notion of minimal pattern occurrences has been signficiantly improved by Burstein \cite{Burstein_Short}.  Since we are introducing the notion of pattern occurrences, we need notation before proceeding.
%
%\begin{notation}
%For a set of permutations containing a pattern, we let $S_n(\tau; r)$ be the set of permutations in $S_n$ containing a pattern $\tau$ exactly $r$ times.
%\end{notation}
%
%To begin, we will state a couple of theorems involving one pattern occurrence of length three.
%
%\begin{theorem}
%\label {Noo_Zeil}
%\cite {Noonan, Noonan_Zeilberger}
%$$|S_n(123;1)| = \frac{3}{n} \binom{2n}{n-3}$$
%\label {Noonan_Thm}
%\end{theorem}
%
%\begin{theorem}
% \cite {Bona_132}
%$$|S_n(231;1)| = \binom{2n-3}{n-3}$$
%\label {Bona_Thm}
%\end{theorem}
%
%There is convenient notation involving generating functions that will allow us to view Theorems \ref {Noonan_Thm} and \ref {Bona_Thm}  in a different light.  We will introduce that notation in the sections stating our results.  The next theorem, involving two pattern occurrences of 123, was conjectured by Noonan and Zeilberger, and later solved by Bona.
%
%\begin{theorem}
%\cite {Noonan_Zeilberger, Bona_132}
%$$|S_n(123;2)| = \frac{59n^2+117n+100}{2n(2n-1)(n+5)} \binom{2n}{n-4}$$
%\end{theorem}
%
%
%\subsection{Block Decomposition}
%
%The method of Block Decomposition was first proposed by Mansour and Vainshtein in \cite{MV_Restricted} and referenced and defined in \cite{MV_Chebyshev} by them as well.  It arose from considering the structure of permutations avoiding 132, as well as permutations with a few occurrences of 132.  Let us look at a straightforward case of 132-avoiding permutations.  Let ${\pi \in S_n(132)}$ and let the entry $n$ occur in the $k$-th position.  Then we have $\pi = (\pi', n, \pi'')$ where, $\pi'$ is a permutation of the set $\{n-k+1, n-k+2, \dots, n-1, n\}$, $\pi''$ is a permutation of the set $[n-k]$, and
% both $\pi'$ and $\pi''$ avoid 132. This representation is called the block decomposition of $\pi$, see Figure \ref{fig:block132_0}.

%\begin{itemize}
%\item $\pi'$ is a permutation of the set $\{n-k+1, n-k+2, \dots, n-1, n\}$,
%\item $\pi''$ is a permutation of the set $[n-k]$, and
%%\item both $\pi'$ and $\pi''$ avoid 132.
%%\end{itemize}
%
%\begin{center}
%\begin{tikzpicture} 
% \foreach \position in {(7,11)}
%    \filldraw [scale = 0.5, color = black] \position circle (8pt);
% \node at (4,5.5) {$n$};
%\draw [line width = 2pt, color=black] (0,2) rectangle (3,5);
% \node at (1.5,3.5){$\pi'$};
%\draw [line width = 2pt, color=black] (4,0) rectangle (6,2);
% \node at (5,1){$\pi''$};
%\draw [line width = 1pt, color=black, dashed] (-0.5,2) -- (6.5,2);
% \node at (7.25,2.1){$n-k$};
%\end{tikzpicture} 
%\begin{figure}[ht]
%\caption{Block decomposition for $\pi \in S_n (132)$}
%\label{fig:block132_0}
%\end{figure}
%\end{center} 
%%
%Since we have introduced the notion of an pattern occurence in the previous section, let us link that notion with that of a block decomposition of exactly one occurence of 132.
%
%\begin{theorem}
%\cite{MV_Chebyshev}
%Let $\pi \in S_n(132;1)$.  Then the block decomposition of $\pi$ can have one of the following forms:
%
%\begin{itemize}
%
%\item there exists k such that $\pi = (\pi', n, \pi'')$ where:
%\begin{itemize}
%\item $\pi'$ is a permutation consisting of the set $\{n-k+1, n-k+2, \dots, n-2, n-1\}$ containing 132 exactly once, and 
%\item $\pi''$ is a permutation consisting of the set $[n-k]$ avoiding 132;
%\end{itemize}
%
%\item there exists k such that $\pi = (\pi', n, \pi'')$ where:
%\begin{itemize}
%\item $\pi'$ is a permutation consisting of the set $\{n-k+1, n-k+2, \dots, n-2, n-1\}$ avoiding 132, and 
%\item $\pi''$ is a permutaiton consisting of the set $[n-k]$ containing 132 exactly once;
%\end{itemize}
%
%\item there exists k, l such that $\pi = (\pi', n-k+1, n, \pi'', n-k+2, \pi''')$ where:
%\begin{itemize}
%\item $\pi'$ is a permutation consisting of the set $\{n-k+3, n-k+4, \dots, n-2, n-1\}$ avoiding 132,  
%\item $\pi''$ is a permutation consisting of the set $\{n-l+1, n-l+2, \dots, n-k-1, n-k\}$ avoiding 132, and
%\item $\pi'''$ is a permutation consisting of the set $[n-k]$ avoiding 132.
%\end{itemize}
%
%\end{itemize}
%
%\end{theorem}
%
%The block decomposition of the first two cases are similar to the earlier block decomposition, see Figure \ref{fig:block132_0}.  The block decomposition of the third case is below, see Figure \ref{fig:block132_1}.
%
%\begin{center}
%\begin{tikzpicture} 
%
% \foreach \position in {(3.5,6.5), (-1.5,5.5), (-0.5, 11.5)}
%    \filldraw [color = black] \position circle (4pt);
%
% \node at (4.75,6.5) {$n-k+2$};
% \node at (-0.25,5.5) {$n-k+1$};
% \node at (0,11.5) {$n$};
%
%\draw [line width = 2pt, color=black] (0,2) rectangle (3,5);
% \node at (1.5,3.5) {$\pi''$};
%
%\draw [line width = 2pt, color=black] (4,0) rectangle (6,2);
% \node at (5,1){$\pi'''$};
%
%\draw [line width = 2pt, color=black] (-6,7) rectangle (-2,11);
% \node at (-4,9){$\pi'$};
%
%\draw [line width = 1pt, color=black, dashed] (-7.5,2) -- (6.5,2);
% \node at (-8.25,2.1){$n-l$};
%
%\end{tikzpicture} 
%\begin{figure}[ht]
%\caption{Block decomposition for $\pi \in S_n (132;1)$}
%\label{fig:block132_1}
%\end{figure}
%\end{center} 
%


% End Section 1 - Preliminaries

%%%%%%%%%%%%%%%%%%%%

% Begin Section 2 - Dumont permutations

\section {Dumont permutations}
%As referenced earlier, there are currently four known types of Dumont permutations.  We will define them after a short introduction.

%\subsection {Introduction}
Dumont permutations are a certain class of permutations shown by Dumont \cite{Dumont} to be counted by the (unsigned) Genocchi numbers \cite[\href{https://oeis.org/A110501}{A110501}]{OEIS}, a multiple of the Bernoulli numbers.  Let the Genocchi and Bernoulli numbers be denoted $G_{2n}$ and $B_{2n}$, respectively, then we have $G_{2n} = 2 (1 - 2^{2n})(-1)^n B_{2n}$.

The exponential generating functions for the unsigned and signed Genocchi numbers are given by
$$\sum_{n=1}^\infty G_{2n} \frac {x^{2n}}{(2n)!} = x \tan \frac {x}{2} \qquad \text{and} \qquad \sum_{n=1}^\infty (-1)^n G_{2n} \frac {x^{2n}}{(2n)!} = \frac {2x}{e^x + 1} - x = - x \tanh \frac {x}{2}.$$

%
%Going forward, all Dumont permutations we consider will be of even length.  So, let $\pi$ be a permutation in $S_{2n}$.

The definition of Dumont permutations of the first and third kinds involves descents and ascents, and thus, the linear structure of permutations. The definition of Dumont permutations of the second and fourth kinds involves fixed points, excedances, and deficiencies, and thus the cyclic structure of permutations. 

\begin{definition}
We say that $i\in[n-1]$ is a \emph{descent} of $\pi$ if $\pi(i) > \pi(i+1)$.  
In this case, we call $\pi(i)$ the \emph{descent top} and $\pi(i+1)$ the \emph{descent bottom}. We also say that $i\in[n-1]$ is an \emph{ascent} of $\pi$ if $\pi(i) < \pi(i+1)$. In this case, we call $\pi(i)$ the \emph{ascent bottom} and $\pi(i+1)$ the \emph{ascent top}.
\end{definition}

\begin{definition}
A \emph{fixed point} (resp. \emph{excedance}, \emph{deficiency}) of a permutation $\pi$ is a position $i$ such that $\pi(i) =  i$ (resp. $\pi(i) > i$, $\pi(i) <  i$).
\end{definition}

We only consider Dumont permutations of even size since it will be clear from the definitions to follow that Dumont permutations of each kind of size $2n+1$ are obtained by appending the fixed point $2n+1$ to the right of a Dumont permutation of the same kind of size $2n$.

\subsection{Dumont permutations of the first and second kinds}

\begin{definition}
A \emph{Dumont permutation of the first kind} (or \emph{Dumont-1 permutation} for short) is a permutation wherein each even entry must be immediately followed by a smaller entry, and each odd entry must be immediately followed by a larger entry, or ends the permutation (i.e. the last entry must be odd).  In other words, for all $i$, $1 \leq i \leq 2n$, and some $k$, $1 \leq k < n$,
\begin{align*}
\pi(i) = 2k &\implies i < 2n \quad \text{and} \quad \pi(i) > \pi(i+1), 
\\
\pi(i) = 2k-1 &\implies \pi(i) < \pi (i+1) \quad \text{or} \quad i = 2n.
\end{align*}
\end{definition}

\begin{example}
\label {Dumont Example 1}
Let $\pi = 435621 \in S_6$. Then $2$, $4$, and $6$, the descent tops of $\pi$, are all even, while the descent bottoms of $\pi$ are $3$ and $5$, and its final entry is $1$, all odd.  Thus, $\pi$ is a Dumont permutation of the first kind.
\end{example}

\begin{definition}
A \emph{Dumont permutation of the second kind } (or \emph{Dumont-2 permutation} for short) is a permutation wherein each entry at an even position is a deficiency, and each entry at an odd position is a fixed point or an excedance.  In other words, for all $i$, $1 \leq i \leq n$,
\[
\pi(2i) < 2i, \qquad
\pi(2i-1) \geq 2i-1.
\]
\end{definition}

\begin{example}
\label {Dumont Example 2}
Let $\rho = 614352 \in S_6$. We see that the deficiencies of $\rho$ occur at positions $2$, $4$ and $6$, all even while excedances occur at positions $1$ and $3$, and a fixed point is at position $5$, all odd. Thus, $\rho $ is a Dumont permutation of the second kind.
\end{example}

\begin{notation} The Dumont permutations of the first (resp. second) kind shall be denoted by $\mathfrak{D}^1_{2n}$ (resp. $\mathfrak{D}^2_{2n}$).
\end{notation}

Thus, from Examples \ref{Dumont Example 1}, and \ref{Dumont Example 2}, we have $\pi \in \mathfrak{D}^1_{6}$, and $\rho \in \mathfrak{D}^2_{8}$.  As stated earlier, Dumont \cite {Dumont} showed that $|\mathfrak{D}^1_{2n}| = |\mathfrak{D}^2_{2n}| = G_{2n+2}$.  There is a simple natural bijection \cite {Dumont_Foata} between Dumont permutations of the first and second kinds, namely, the \emph{Foata's fundamental transformation} \cite{Foata} that maps left-to-right maxima to cycle maxima.

%The reader is invited to see \cite{Jones} for details.

\begin{definition}
An entry $\pi(i)$ of a permutation $\pi$ is a \emph{left-to-right maximum} if $\pi(i)>\pi(j)$ for all $j<i$.
%(and left-to-right minimum as in Definition \ref {LR Minimum}).
\end{definition}

In particular, the first entry of $\pi$ is always a left-to-right maximum $\pi$.

The map $f$, called the \emph{Foata's fundamental transformation} \cite{Foata}, is defined as follows:

\begin{itemize}
\item Start with a Dumont permutation of the first kind, say $\pi$.
\item Insert parentheses to display the permutation in cycle notation, so that each cycle starts with a left-to-right maximum.
\end{itemize}

Then $f$ is a bijection, and $\pi\in\mathfrak{D}^1_{2n}$ if and only if $f(\pi)\in\mathfrak{D}^2_{2n}$.

\begin{example}
For $\pi = 435621 \in \mathfrak{D}^1_{6}$, we have the left-to-right maxima $\pi(1)=4$, $\pi(3)=5$, and $\pi(4)=6$, so
$$f(\pi)= (43) (5) (621) \, = \, 614352.$$
Notice that $f(\pi) = 614352 \in \mathfrak{D}^2_{6}.$
\end{example}

%Now let us describe the inverse $g=f^{-1}$ of $f$.
%\begin{itemize}
%\item Start with a Dumont permutation of the second kind, say $\rho$.
%\item Rewrite the permutation implementing cycle notation.
%\item Within each cycle, write the maximal element first.
%\item If need be, reorder the cycles in increasing order of maximal elements.  Note, the beginning of each cycle will create the left-to-right maxima needed.
%\item Erase the parentheses, to display the permutation in one-line form.
%\item This new permutation, $g(\rho)$, is a Dumont permutation of the first kind.
%\end{itemize}
%
%\begin{example}
%Let $\rho = 614352 \in \mathfrak{D}^2_{6}$.  We have,
%$$614352 \, = \, (162) (34) (5) \, = \, (621) (43) (5) \, = \, (43) (5) (621) \, \mapsto \, 435621.$$
%It follows that $g(\rho) = 435621 \in \mathfrak{D}^1_{6}.$
%And clearly, $g(f(\pi)) = \pi$ and $f(g(\rho)) = \rho$.
%\end{example}



\subsection{Dumont permutations of the third and fourth kinds}

With respect to the next two types of Dumont permutations, Kitaev and Remmel \cite {Kitaev_Remmel_1, Kitaev_Remmel_2} first conjectured that sets of permutations where each descent is from an even value to an even value are also counted by the Genocchi numbers.  Burstein and Stromquist \cite {Burstein_Third} proved this conjecture and called those sets of permutations Dumont permutations of the third kind. As with Dumont-1 and Dumont-2 permutations, the Foata's fundamental transformation maps Dumont permutations of the third kind onto a related equinumerous set that Burstein and Stromquist  \cite {Burstein_Third} called Dumont permutations of the fourth kind.  We now introduce these two sets.

\begin{definition}
A \emph{Dumont permutation of the third kind} (or \emph{Dumont-3 permutation} for short) of size $2n$ ($n\in\mathbb{N}$) is a permutation where each descent is from an even value to an even value, that is all descent tops and all descent bottoms are even.  In other words, for all $i$, $1 \leq i \leq 2n-1$,
$$\pi(i) > \pi(i+1) \implies \pi (i) = 2l \, \text{ and }\, \pi(i+1) = 2k$$
for some $k$ and $l$, $1 \leq k < l \leq n$.
\end{definition}

\begin{example}
\label {Dumont Example 3}
Let $\pi = 16238457 \in S_8$. Then the descent tops of $\pi$ are $6$ and $8$ are and the descent bottoms are $2$ and $4$, all even entries. Thus, $\pi$ is a Dumont permutation of the third kind.
\end{example}

\begin{definition}
A \emph{Dumont permutation of the fourth kind} (or \emph{Dumont-4 permutation} for short) of size $2n$ ($n\in\mathbb{N}$) is a permutation where deficiencies must be even values at even positions. In other words, for all $i$, $1 \leq i \leq 2n$,
$$\pi(i) < i \implies i = 2u \, \text{ and } \, \pi(i) = 2v$$
for some $1 \leq v < u \leq n$.
\end{definition}

%Note that, unlike in Dumont-3 permutations, where $k$ must be less than $l$, in Dumont-4 permutations there is no correlation between $u$ and $v$.

\begin{example}
\label {Dumont Example 4}
Let $\rho = 13657284 \in S_8$. Then the deficiencies of $\rho$ are $\pi(6)=2$, $\pi(8)=4$, where all the positions and entries involved are even.  Thus, $\rho$ is a Dumont permutation of the fourth kind.
\end{example}

\begin{notation} Following the same notation as earlier, the Dumont permutations of the third (resp. fourth) kind shall be denoted by $\mathfrak{D}^3_{2n}$ (resp. $\mathfrak{D}^4_{2n}$).
\end{notation}

Thus, from Examples \ref {Dumont Example 3}, and \ref {Dumont Example 4}, we have $\pi \in \mathfrak{D}^3_{8}$, and $\rho \in \mathfrak{D}^4_{8}$.  The same bijection, Foata's fundamental transformation, yields $|\mathfrak{D}^3_{2n}| = |\mathfrak{D}^4_{2n}| = G_{2n+2}$. (In fact, notice that $f(16238457)=13657284$ in Examples \ref{Dumont Example 3} and \ref{Dumont Example 4} above.)


We now reference the proof \cite {Burstein_Third} that  $|\mathfrak{D}^1_{2n}| = |\mathfrak{D}^3_{2n}| = G_{2n+2}$, thus showing that all four kinds of Dumont permutations are counted by the Genocchi numbers.

%%Also keeping with Example \ref {Dumont Example 3}, we will use the same simple natural bijection to show that $|\mathfrak{D}^3_{2n}| = |\mathfrak{D}^4_{2n}| = |G_{2n+2}|$.
%
%\begin{example}
%$$\mathfrak{D}^3_{8} \ni 16238457 \xmapsto{\textit{f}} (1) (623) (8457) = 13657284 \in \mathfrak{D}^4_{8}$$
%$$\mathfrak{D}^4_{8} \ni 13657284 = (1) (236) (4578) = (1) (623) (8457) \xmapsto{\textit{g}} 16238457 \in \mathfrak{D}^3_{8}$$
%\end{example}



%Again, we have $g(f(\pi)) = \pi$ and $f(g(\rho)) = \rho$.  And since we have shown that  $|\mathfrak{D}^1_{2n}| = |\mathfrak{D}^2_{2n}| = |G_{2(n+1)}|$ and  $|\mathfrak{D}^3_{2n}| = |\mathfrak{D}^4_{2n}| = |G_{2(n+1)}|$, we will refererence the proof that  $|\mathfrak{D}^1_{2n}| = |\mathfrak{D}^3_{2n}| = |G_{2(n+1)}|$ from Burstein and Stromquist \cite {Burstein_Third}, thus showing that all four kinds of Dumont permutations are counted by the Genocchi numbers.

%
%\subsection{Dumont-1 and Dumont-3 permutations are equinumerous}
%
%Burstein and Stromquist \cite {Burstein_Third} showed that the strategy for proving that the sets of Dumont-1 and Dumont-3 permutations have the same cardinality lies in using a combinatorial structure called a \emph{signature}, counting the Dumont-1 and Dumont-3 permutations corresponding to each possible signature, via the \emph{signature function} for each signature, and then showing that the values of the associated signature functions are equal in the corresponding Dumont-1 and Dumont-3 permutations.  Since, for any given signature, there are as many Dumont-3 permutations as Dumont-1 permutations with that signature, Burstein and Stromquist \cite {Burstein_Third} conclude that Dumont-1 and Dumont-3 permutations are equinumerous.
%
%\begin{definition} 
%\cite {Burstein_Third}
%A \emph{signature} of order $k$ and size $2n$ consists of the following:
%
%\begin{itemize}
%\item a set of $k$ even numbers from $[2n]$ called \emph{peaks}, and 
%\item a set of $k$ odd numbers from $[2n]$ called \emph{pits}.
%\end {itemize}
%
%\noindent Given a particular signature, let:
%\begin{itemize}
%\item $\pk(x)$ be the number of peaks that are at least $x$, and 
%\item $\pt(x)$ be the number of pits that are at least $x$.
%\end {itemize}
%
%\noindent The associated \emph{signature function} has domain $[2n]$ and is given by,
%
%$$
%f(x) = \left\{
%        \begin{array}{ll}
%            \pk(x) - \pt(x) & \quad \text{if $x$ is even}\\
%            \pk(x) - \pt(x) + 1 & \quad \text{if $x$ is odd.}
%        \end{array}
%    \right.
%$$
%\end{definition}
%
%\vspace {0.1in}
%
%\begin{example}
%$(\{6, 8\}, \{1, 3\})$ is a signature of size $8$ and order $2$.  Its signature function is given by 
%
%\begin{center}
%\begin{tabular}{c||c|c|c|c|c|c|c|c}
%	$x$ & 1 & 2 & 3 & 4 & 5 & 6 & 7 & 8\\
%	\hline 
%	$f(x)$ & 1 & 1 & 2 & 2 & 3 & 2 & 2 & 1
%\end{tabular}
%\end{center}
%\vspace{0.1in}
%and the product of the values is $48$. 
%\end{example}
%
%The specific structure of a signature of Dumont-1 permutations is as follows:
%
%\begin{definition} 
%\cite {Burstein_Third}
%For signatures of Dumont-1 permutations, the \emph{peaks} are the even values that do not end descents (i.e. are not descent bottoms), and the \emph{pits} are the odd values that do end descents (i.e. are  descent bottoms.)  The \emph{signature} of a Dumont-1 permutation consists of the set of its peaks and the set of its pits, generally in numeric order.
%\end{definition} 
%
%\begin{example}
%Let $\pi = 58364217 \in \mathfrak{D}^1_8$.  Its peaks are 6 and 8, its pits are 1 and 3, and therefore its signature as $(\{6,8\}, \{1,3\}).$
%\end{example}
%
%Now the specific structure of a signature of Dumont-3 permutations is as follows:
%
%\begin{definition} 
%\cite {Burstein_Third}
%For signatures of Dumont-3 permutations, the \emph{peaks} are the even values that do not end descents (i.e. are not descent bottoms), and the \emph{pits} are the even values that do begin descents (i.e. are descent tops.)  A peak that is also a pit is called a \emph{singleton}.  The \emph{signature} of a Dumont-3 permutation consists of 
%\begin{itemize}
%\item a list of its peaks, and
%\item a list of the values of the form $x-1$ for each pit $x$.
%\end{itemize}
%The second list contains odd numbers that we call \emph{pit list entries}.
%\end{definition} 
%
%\begin{example}
%Let $\pi = 15846237 \in \mathfrak{D}^3_8$.  Its peaks are 6 and 8, its pits are 2 and 4, so the pit list entries are 1 and 3, and therefore its signature is $(\{6,8\}, \{1,3\})$.
%\end{example}
%
%Lastly, we will now state the main theorem of this section and its corollaries without proof.
%
%\begin{theorem}
%\cite {Burstein_Third}
%Given a signature of size $2n$ and its associated signature function $f$, the number of Dumont-1 permutations with that signature and the number of Dumont-3 permutations with that signature are both exactly $\prod_{x=1}^{2n} f(x)$.
%\end{theorem} 
%
%\begin{corollary} 
%\cite {Burstein_Third}
%Given a signature, there are as many Dumont-3 permutations with that signature as Dumont-1 permutations with that signature.
%\end{corollary} 
%
%\begin{corollary} 
%\cite {Burstein_Third}
%For each $n$, there are as many Dumont-3 permutations of size $2n$ as Dumont-1 permutations of size $2n$. That is, $|\mathfrak{D}^3_{2n}| = |\mathfrak{D}^1_{2n}|.$
%\end{corollary} 
%
%We now have our desired result that all four kinds of Dumont permutations are counted by the Genocchi numbers.
%

\subsection{Pattern avoidance in Dumont-1 and Dumont-2 permutations}

There have been several enumerations of pattern-restricted Dumont-1 and Dumont-2 permutations, mostly by Mansour, Burstein, Elizalde, and Ofodile \cite{Mansour, Burstein_Restricted, Burstein_Dyck, Ofodile}.  We will first reproduce their theorems for enumerations of Dumont-1 permutations avoiding patterns of length three.  


\subsubsection{Patterns of length three}

For the six patterns of length three, all the corresponding pattern-restricted sets have been enumerated, both in Dumont-1 and Dumont-2 permutations.

\begin{theorem}[\cite{Mansour}] 
For all $n\ge 0$,
\[
|\mathfrak{D}^1_{2n} (132)| = |\mathfrak{D}^1_{2n} (231)| = |\mathfrak{D}^1_{2n} (312)| = C_n = \frac{1}{n+1} \binom{2n}{n}.
\]
\end{theorem}

\begin{theorem}[\cite{Burstein_Restricted}] 
For all $n\ge 1$,
\[
|\mathfrak{D}^1_{2n} (213)| = 
%\sum_{k=1}^{n-1} C_{n-1-k}C_{k-1} = 
C_{n-1}.
\]
\end{theorem}

It is important to note here that although reverses, complements, and inverses of patterns create symmetry classes and Wilf-equivalences in the set of all permutations, Dumont permutations are not closed under any symmetry operations, so that applying a symmetry operation to a pattern does not necessarily yield a pattern-avoiding set of the same cardinality. For example, patterns $132$ and $312$ are complements of each other, and $|\mathfrak{D}^1_{2n} (132)| = |\mathfrak{D}^1_{2n} (312)| = C_n$; on the other hand, patterns $213$ and $231$ are also complements of each other, but $|\mathfrak{D}^1_{2n} (213)| = C_{n-1} \ne C_n =  |\mathfrak{D}^1_{2n} (231)|$. 

\begin{theorem}[\cite{Burstein_Restricted}] 
For all $n\ge 0$,
\[
|\mathfrak{D}^1_{2n} (321)| = 1, \quad \text{namely,} \quad \mathfrak{D}^1_{2n} (321) = \{2, 1, 4, 3, \dots, 2n, 2n-1\}.
\]
\end{theorem}

\begin{theorem}[\cite{Burstein_Restricted}] 
For all $n\ge 3$,
\[
|\mathfrak{D}^1_{2n} (123)| = 4, \quad \text{namely,} \quad \mathfrak{D}^1_{2n} (123) = \{(2n-1, 2n, 2n-3, 2n-2, \dots, 7, 8, \pi)\, |\, \pi \in \mathfrak{D}^1_{6} (123)\},
\]
where
\[
\mathfrak{D}^1_{6} (123) = \{436215, 562143, 563421, 564213\}.
\]
\end{theorem}

Now, with respect to enumerations of Dumont-2 permutations avoiding patterns of length three, we can start with the makeup of a Dumont-2 permutation.  It is straightforward that for all $n \ge 3$, 
$$|\mathfrak{D}^2_{2n} (123)| = |\mathfrak{D}^2_{2n} (132)| = |\mathfrak{D}^2_{2n} (213)| = 0.$$  
This follows from the fact that every even position is a deficiency, every odd position is an excedance or a fixed point, and for any Dumont-2 permutation $\pi(2) = 1$ and $\pi(2n-1) = 2n$ or $2n-1$,  which eventually leads to the conclusion that it is impossible to avoid 123, 132, and 213.

\begin{theorem}[\cite{Burstein_Restricted}]
For all $n\ge 1$,
\[
|\mathfrak{D}^2_{2n} (231)| = 2^{n-1}.
\] 
\end{theorem}

\begin{theorem}[\cite{Burstein_Restricted}]
 For all $n\ge 0$,
\[
|\mathfrak{D}^2_{2n} (312)| = 1, \quad \text{namely} \quad \mathfrak{D}^2_{2n} (312) = \{2, 1, 4, 3, \dots, 2n, 2n-1\}.
\] 
\end{theorem}

\begin{theorem}[\cite{Mansour}]
For all $n\ge 0$,
\[
|\mathfrak{D}^2_{2n} (321)| = C_n = \frac{1}{n+1} \binom{2n}{n}.
\]
\end{theorem}


\subsubsection{Patterns of length four}
\label{length four_thms}

Now, for Dumont-1 and Dumont-2 permutations avoiding patterns of length four, there are several cases which are open, as well as several enumerations of permutations that avoid two patterns of length four simultaneously.  We begin with Dumont-2 permutations avoiding a single pattern of length four.  
%Later in our paper, we use methods similar to that used in the first of these results below, and thus its proof will be reconstructed here.

\begin{theorem}[\cite{Burstein_Restricted}]
\label{3142_thm}
For all $n\ge 0$,
\[
|\mathfrak{D}^2_{2n} (3142)| = C_{n}.
\]
\end{theorem}
%
%\begin{proof}
%Let $\pi \in \mathfrak{D}^2_{2n} (3142).$  To begin, since $\pi$ is a Dumont-2 permutation, $\pi(2) = 1.$  We will consider two cases based on parity of the first entry, $\pi(1)$.
%
%Suppose that $\pi(1)$ is odd, say $\pi(i)=2k-1$ for some $2 \le k \le n$ (we cannot have $k=1$ since $\pi(2)=1$).  Then $\pi(2k-1)$, in an odd position, cannot be a fixed a point, and thus must be strictly greater than $2k-1$.  Now, if there exists an entry $\pi(i) < 2k-1$ to the right of $\pi(2k-1)$, then the permutation $$(\pi(1), \pi(2), \pi(2k-1), \pi(i))$$ is an occurrence of the pattern 3142 in $\pi$.  Therefore, all $\pi(i) \in [2,2k-2]$ must occur at positions $i \in [3, 2k-2]$, which is impossible by the pigeon hole principle.  Therefore, $\pi(1)$ must be even, in other words $\pi(1) = 2k$ for some $k \in [1,n]$.
%
%Now, since $\pi(2k+1) \ge 2k$, as before we have $2 \le \pi(i) \le 2k-1$ only if $3 \le i \le 2k.$  That is all $\pi(i) \in [2,2k-1]$ must occur at positions $i \in [3, 2k]$. But the number of these values is the same as the the number of the positions they can possibly occupy, so $\pi(i) \in [2k+1,2n]$ must occur at positions $i \in [2k+1,2n]$. In other words, we have $$\pi = (2k, 1, \pi_1 + 1, \pi_2 + 2k),$$ where
%\begin{itemize}
%\item $\pi_1$ is a permutation of the set $[2k-2]$ avoiding 3142,
%\item $\pi_1 + 1$ is a block in $\pi \in \mathfrak{D}^2_{2n} (3142)$ of length $2k-2$ starting at position $3$, and
%\item $\pi_2 \in \mathfrak{D}^2_{2n-2k} (3142).$
%\end{itemize}
%We also see that the composition of the reverse and complement of $\pi_1$, we denote it $\pi'$, belongs to the set of Dumont-2 permutations of length $2k-2$ that avoid 3142.  That is $\pi' = (\pi_1)^{rc} \in \mathfrak{D}^2_{2k-2} (3142)$ since $(3142)^{rc} = 3142.$  Thus $\pi \in \mathfrak{D}^2_{2n} (3142)$ if and only if $$\pi = (2k, 1, (\pi')^{rc} + 1, \pi'' + 2k),$$ for some $k \in [1,n]$, and any $\pi' \in \mathfrak{D}^2_{2k-2} (3142)$ and $\pi'' \in \mathfrak{D}^2_{2n-2k} (3142)$.  
%
%Lastly, if we let $a_n = |\mathfrak{D}^2_{2n} (3142)|$, then $a_0 = 1$ and $\sum^n_{k=1} a_{k-1} a_{n-k}$ for all $n \ge 1$, thus $a_n = C_n.$
%\end{proof}

For pattern $4132$, note that $321$ is a subsequence of $4132$, therefore $\mathfrak{D}^2_{2n} (321) \subseteq \mathfrak{D}^2_{2n} (4132)$. The following theorem shows that, in fact, the two sets are equal.

\begin{theorem}[\cite{Burstein_Dyck}]
For all $n\ge 0$,
\[
|\mathfrak{D}^2_{2n} (4132)| = |\mathfrak{D}^2_{2n} (321)| = C_n.
\]
\end{theorem}

\begin{theorem}[\cite{Burstein_Dyck}]
For all $n\ge 0$,
\[
|\mathfrak{D}^2_{2n} (2143)| = a_n a_{n+1},
\]
where
\[
a_{2m} = \frac{1}{2m+1} \binom{3m}{m} \qquad \text{and} \qquad a_{2m+1} = \frac{1}{m+1} \binom{3m+1}{m}.
\]
%and
%\[
%a_{2m+1} = \frac{1}{2m+1} \binom{3m+1}{m+1} = \frac{1}{m+1} \binom{3m+1}{m}.
%\]
\end{theorem}

\vspace{0.2in}

To date, there are no other enumerations of Dumont-1 and Dumont-2 permutations avoiding a single pattern of length four. However, we conjecture a Wilf-equivalence on Dumont-1 permutations of patterns $2143$ and $3421$ (see Section~\ref {B-J Conjecture}).

Now, we will consider simultaneous avoidance of a pair of patterns. We begin with Dumont-1 permutations.

\begin{theorem}[\cite{Burstein_Dyck}]
For all $n\ge 0$,
\[
|\mathfrak{D}^1_{2n} (1342, 1423)| = |\mathfrak{D}^1_{2n} (2341, 2413)| = |\mathfrak{D}^1_{2n} (1342, 2413)| =s_{n+1},
\]
where $s_n$ is the $n$-th little Schr\"oder number (\cite[A001003]{OEIS}).
\end{theorem}

Note that $(s_n)=(1, 1, 3, 11, 45, 197, 903, \dots)$ is recursively given by $s_1 = 1$ and
\[
s_{n+1} = -s_n + 2 \sum_{k=1}^n s_k s_{n-k}, \quad n\ge 2,
\]
with the generating function
\[
s(x) = \sum_{n=1}^\infty s_n x^n = \frac{1+x-\sqrt{1-6x+x^2}}{4}.
\]
%
%\begin{proof}
%Let $\pi \in \mathfrak{D}^1_{2n} (1342, 1423).$  Then $\pi$ can be one of the two following cases:
%\begin{itemize}
%\item $(\pi_1, 2n-1, 2n, \pi_2)$ where $\pi_1$ is empty, or $\pi_1 > \pi_2$ where neither are empty.
%\item $(\pi_1, 2n, \pi_2, 2n-1)$ where $\pi_1 > \pi_2$, or $\pi_1$ is empty, or $\pi_2$ is empty. 
%\end{itemize}
%
%In the first case $\pi$ avoids $1342$, and in the second case $\pi$ avoids $1423$.  In either case, if $\pi_1$ is nonempty, then its minimum entry is odd.  Therefore, the maximum entry of $\pi_2$ is even, regardless of whether or not $\pi_1$ is nonempty.  Now, for $\pi$ of either case, the entries $2n-1$ and $2n$ cannot be part of any instance of $1342$ or $1423$.  Thus, $\pi$ avoids both $1342$ and $1423$ if and only if both $\pi_1$ and $\pi_2$ avoid both $1342$ and $1423$.  Hence, $\pi \in \mathfrak{D}^1_{2n} (1342, 1423)$ if and only if,
%\begin{itemize}
%\item $(\pi' + 2k, 2n-1, 2n, \pi'')$ for $1 \le k \le n-1$ where:
%\begin{itemize}
%\item  $\pi' \in \mathfrak{D}^1_{2n-2k-2} (1342, 1423),$
%\item  $\pi'' \in \mathfrak{D}^1_{2k} (1342, 1423)$; or,
%\end{itemize}
%\item $(\pi' + 2k, 2n, \pi'', 2n-1)$ for $0 \le k \le n-1$ where again:
%\begin{itemize}
%\item  $\pi' \in \mathfrak{D}^1_{2n-2k-2} (1342, 1423),$
%\item  $\pi'' \in \mathfrak{D}^1_{2k} (1342, 1423).$
%\end{itemize}
%\end{itemize}
%
%\noindent Therefore, for $a_n = |\mathfrak{D}^1_{2n} (1342, 1423)|$, we have $a_0 = 1$ and, 
%$$a_n = -a_{n-1} + 2 \sum_{k=0}^{n-1} a_k \, a_{n-1-k},$$ so $\{a_n\}_{n \ge 0}$ satisifies the same recurrence relation as $\{s_{n+1}\}_{n \ge 0}$, thus $a_n = s_{n+1}$.
%\end{proof}
%
%The next two enumerations of Dumont permutations avoiding a pair of patterns use the same methodology with the exact same results.
%
%\begin{theorem}
%\cite{Burstein_Dyck} For all $n\ge 0$,
%$$|\mathfrak{D}^1_{2n} (2341, 2413)| =s_{n+1}.$$ 
%\end{theorem}
%
%\begin{theorem}
%\cite{Burstein_Dyck} For all $n\ge 0$,
%$$|\mathfrak{D}^1_{2n} (1342, 2413)| =s_{n+1}.$$ 
%%\end{theorem}
%
%The next two enumerations of Dumont permutations avoiding a pair of patterns have similar proofs that have been introduced, thus only the theorems will be restated.

\begin{theorem}[\cite{Burstein_Dyck}] 
For all $n\ge 1$,
\[
|\mathfrak{D}^1_{2n} (231, 4213)| = 1, \quad \text{namely,} \quad \mathfrak{D}^1_{2n} (231, 4213) = \{(2, 1, 4, 3, \dots 2n, 2n-1)\}.
\]
\end{theorem}

\begin{theorem}[\cite{Burstein_Dyck}] 
For all $n\ge 1$, $|\mathfrak{D}^1_{2n} (1342, 4213)| = 2^{n-1}$.
\end{theorem}

\begin{theorem}[\cite{Burstein_Dyck}] 
For all $n\ge 3$,
$|\mathfrak{D}^1_{2n} (2341, 1423)| = b_n$, where $b_n$ satisfies the recurrence relation 
\[
b_n = 3b_{n-1} + 2b_{n-2} \quad \text{for} \quad n \ge 3, \quad \text{with} \quad b_0 = 1, b_1 = 1, b_2 = 3.
\]
\end{theorem}

Note that the sequence $(b_n)=(1, 1, 3, 11, 39, 139, 495, \dots)$ is \cite[A007482]{OEIS} shifted one position to the right. In other words, $\mathfrak{D}^1_{2n} (2341, 1423)$ is equinumerous to the set of of subsets of $[2n-2]$ where each odd element $m$ has an even neighbor ($m-1$ or $m+1$).

% End Section 2 - Dumont permutations

%%%%%%%%%%%%%%%%%%%%

% Begin Section 3 - Avoidance on Dumont permutations of the fourth kind

\section{Avoidance on Dumont permutations of the fourth kind}

In this section, we will consider pattern avoidance on Dumont-4 permutations. In previous pattern-avoidance literature, the first nontrivial cases to be analyzed were patterns of length $3$. Since $1$ cannot be a deficiency in a Dumont-4 permutation, it must be a fixed point, so all Dumont-4 permutations start with $1$. Thus, we will also consider avoiding patterns $\pi= (1,\pi'+1)$, where $\pi'$ is a permutation in $S_3$, i.e. $\pi\in\{1234, 1243, 1324, 1342, 1423, 1432\}$. Note that in all but the first two cases, i.e. if $\pi'$ does not start with $1$, we have 
\[
\mathfrak{D}^4_{2n}(\pi)=\mathfrak{D}^4_{2n}(\pi').
\]

\subsection{Enumerating Dumont-4 permutations avoiding Dumont-4 permutations of length four}

We will begin by considering Dumont-4 permutations avoiding patterns of length $4$, which are themselves Dumont-4 permutations, that is $\pi\in\{1234, 1342, 1432\}$.

Then in the next section we will consider the remaining cases, where $\pi\in\{1243, 1324, 1423\}$.

\subsubsection{$\mathfrak{D}^4_{2n}(1234)$}

\begin{theorem}
$|\mathfrak{D}^4_{2n}(1234)| = 0$ \ for $n \ge 4$.
\end{theorem}

Note that $|\mathfrak{D}^4_{2n}(1234)| = 1, 1, 2, 4$, for $n=0, 1, 2, 3$, respectively. The Dumont-4 permutations of length at most 6 avoiding pattern 1234 are $\epsilon$, $12$, $1342$, $1432$, $132654$, $136254$, $143265$, $143652$, where $\epsilon$ denotes the empty permutation.

\begin{proof}
Recall that in Dumont-4 permutations, $\pi(1)=1$, and $\pi(2n-1)=2n-1$ or $2n$. In addition to that, for $n \ge 2$, we have $\pi(2)=3$ or $\pi(3)=3$. 

Since $\pi(1)=1$, the three conditions above mean that, of the $2n-5$ entries in $[4, 2n-2]$, at most one is to the left of $3$ and at most one is to the right of $\pi(2n-1)$. This leaves at least $2n-7\ge 1$ entries in $[4,2n-2]$ that are to the right of $3$ and to the left of $\pi(2n-1)$. Any such entry, together with $1$, $3$, and $\pi(2n-1)$ would form an occurrence of pattern $1234$.
%
%From these three properties, the following four cases follow for any $\pi \in \mathfrak{D}^4_{2n}, n \ge 4$:
%\begin{enumerate} 
%\item $\pi = (1, 3, \underline{\hspace{0.3cm}}, \underline{\hspace{0.3cm}}, \underline{\hspace{0.3cm}}, \hspace{0.05cm} . . .  \hspace{0.07cm} ,  \underline{\hspace{0.3cm}}, 2n, \underline{\hspace{0.3cm}})$,
%\item $\pi = (1, 3, \underline{\hspace{0.3cm}}, \underline{\hspace{0.3cm}}, \underline{\hspace{0.3cm}}, \hspace{0.05cm} . . .  \hspace{0.07cm} ,  \underline{\hspace{0.3cm}}, 2n-1, \underline{\hspace{0.3cm}})$,
%\item $\pi = (1, \underline{\hspace{0.3cm}}, 3, \underline{\hspace{0.3cm}}, \underline{\hspace{0.3cm}}, \hspace{0.05cm} . . .  \hspace{0.07cm} ,  \underline{\hspace{0.3cm}}, 2n, \underline{\hspace{0.3cm}})$,
%\item $\pi = (1, \underline{\hspace{0.3cm}}, 3, \underline{\hspace{0.3cm}}, \underline{\hspace{0.3cm}}, \hspace{0.05cm} . . .  \hspace{0.07cm} ,  \underline{\hspace{0.3cm}}, 2n-1, \underline{\hspace{0.3cm}})$,
%\end{enumerate}
%where the underscores denote the positions to be filled later.
%
%We will now show that in each of these four cases, it is impossible to avoid 1234.
%
%\bigskip
%
%\underline {Case 1}: $\pi = (1, 3, \underline{\hspace{0.3cm}}, \underline{\hspace{0.3cm}}, \underline{\hspace{0.3cm}}, \hspace{0.05cm} . . .  \hspace{0.07cm} ,  \underline{\hspace{0.3cm}}, 2n, \underline{\hspace{0.3cm}})$.\\
%If $2n-1$ preceded $2n$, we would have a 1234-occurrence $(1,3,2n-1,2n)$, thus $2n-1$ follows $2n$, i.e. $2n-1$ is last. So now we have $\pi = (1, 3, \underline{\hspace{0.3cm}}, \underline{\hspace{0.3cm}}, \underline{\hspace{0.3cm}}, \hspace{0.05cm} . . .  \hspace{0.07cm} ,  \underline{\hspace{0.3cm}}, 2n, 2n-1)$.  Thus, all remaining entries occur between 3 and $2n$, any of which (excluding the entry 2) would yield a 1234-occurrence.
%\bigskip
%
%\underline {Case 2}: $\pi = (1, 3, \underline{\hspace{0.3cm}}, \underline{\hspace{0.3cm}}, \underline{\hspace{0.3cm}}, \hspace{0.05cm} . . .  \hspace{0.07cm} ,  \underline{\hspace{0.3cm}}, 2n-1, \underline{\hspace{0.3cm}})$.\\
%If the entry $2n$ occurred to the right of the entry $2n-1$, that is if $2n$ were last, we would have a 1234-occurrence. Thus, $2n$ occurs to the left of the entry $\pi(2n-1)=2n-1$. Therefore, $\pi(2n)\le 2n-2$. Note that only one entry is greater than $2n-1$ and only one other entry occurs to the right of value $2n-1$. That leaves $2n-6\ge 2$ values (i.e. all except 1, 2, 3, $\pi(2n)$, $2n-1$, $2n$) that are above and the right of 3 and below and to the left of $2n-1$. Any such value, together with entries 1, 3, and $2n-1$ forms an occurrence of pattern 1234.
%
%\bigskip
%
%\underline {Case 3}: $\pi = (1, \underline{\hspace{0.3cm}}, 3, \underline{\hspace{0.3cm}}, \underline{\hspace{0.3cm}}, \hspace{0.05cm} . . .  \hspace{0.07cm} ,  \underline{\hspace{0.3cm}}, 2n, \underline{\hspace{0.3cm}})$.\\
%To avoid 1234, the entry $2$ could not occur in the second position, and the entry $2n-1$ would have to occur in the second or last position. WLOG, consider the entry $2n-2$, and let it occur last or second. In other words, let $2n-2$ occur in the position where $2n-1$ is not.  So we have either $\pi = (1, 2n-2, 3, \underline{\hspace{0.3cm}}, \underline{\hspace{0.3cm}}, \hspace{0.05cm} . . .  \hspace{0.07cm} ,  \underline{\hspace{0.3cm}}, 2n, 2n-1)$, or $\pi = (1, 2n-1, 3, \underline{\hspace{0.3cm}}, \underline{\hspace{0.3cm}}, \hspace{0.05cm} . . .  \hspace{0.07cm} ,  \underline{\hspace{0.3cm}}, 2n, 2n-2)$.  As in the previous cases, all remaining entries (excluding the entry 2) would yield a 1234-occurrence.
%
%\bigskip
%
%\underline {Case 4}: $\pi = (1, \underline{\hspace{0.3cm}}, 3, \underline{\hspace{0.3cm}}, \underline{\hspace{0.3cm}}, \hspace{0.05cm} . . .  \hspace{0.07cm} ,  \underline{\hspace{0.3cm}}, 2n-1, \underline{\hspace{0.3cm}})$.\\
%Again, to avoid 1234, the entry $2$ could not occur in the second position, and the entry $2n$ could not occur in the last position.  So let the entries $2$ and $2n$ occur between positions $3$ and $2n-1$, say  $\pi = (1, \underline{\hspace{0.3cm}}, 3, 2, \underline{\hspace{0.3cm}}, \underline{\hspace{0.3cm}}, \hspace{0.05cm} . . .  \hspace{0.07cm} ,  \underline{\hspace{0.3cm}}, 2n, 2n-1, \underline{\hspace{0.3cm}})$. To avoid 1234, all of the remaining entries $(4, 5, \hspace{0.05cm} . . .  \hspace{0.07cm} , 2n-2)$ would have to be second or last. Since $n \ge 4$, there are at least three such entries vying for two positions, thus yielding a 1234-occurrence.
\end{proof}


\subsubsection{$\mathfrak{D}^4_{2n}(1342)$}

As noted earlier, since the entry following $1$ is not $2$, it follows that $\mathfrak{D}^4_{2n}(1342)=\mathfrak{D}^4_{2n}(231)$.

\begin{theorem}
\label{1342_theorem}
$|\mathfrak{D}^4_{2n}(1342)| = 2^{n-1}$, \text{for} $n \ge 1$.
\end{theorem}

We will first prove that all odd entries are fixed points.

\begin{lemma}
\label{lm:odd_fixed}
Let  $\pi \in \mathfrak{D}^4_{2n}(1342)$.  Then $\pi(2k-1) = 2k-1$ for all $k$, $1 \le k \le n$.
\end{lemma}

\begin{proof}
By definition of a Dumont-4 permutation $\pi(1)=1$, in other words, $1$ is a fixed point. 

%Now consider $\pi(3)$, and suppose $\pi(3) \ne 3$.  The only possibilities for an odd entry are to be a fixed point or an excedance. Since $\pi(3)$ is not a fixed point and $\pi(1)=1$, then $\pi(2)=3$.  Considering the entry 2, it cannot fall in the third position, if so it would be a deficiency but not in an even position. Thus, $\pi (3) \ge 4$, and the entry 2 must occur to the right of $\pi(3)$, yielding a 1342-occurrence.  Therefore $\pi(3)$ must be equal to 3, that is 3 is fixed.

Now assume that all odd entries from $1$ through $2j-1$ are fixed points.  Consider the entry in the next odd position, $\pi(2j+1)$, and as before, suppose $\pi(2j+1) \ne 2j+1$.  It follows that 
\begin{itemize}
\item the entry $2j+1$ must be to the left of $\pi(2j+1)$, 
\item $\pi(2j+1)>2j+1$ as an odd position cannot be a deficiency, and 
\item at least one even entry $2l\le 2j$ must be to the right of $\pi(2j+1)$. This is because all odd entries at most $2j+1$ are to the left of $\pi(2j+1)$, so if all even entries less than $2j+1$ are also to the left of $2j+1$, then there will be $2j+1$ entries to the left of position $2j+1$, which is impossible.
\end{itemize}

This yields a $1342$-occurrence $(1, 2j+1, \pi(2j+1), 2l)$.  Therefore $\pi(2j+1)=2j+1$, that is $2j+1$ is fixed.

By induction, the lemma is proved.
\end{proof}

Next, we will prove that if we have a deficiency at position $2k$, then it must be the entry $2k-2$.


\begin{lemma}
\label{lm:def=2k-2}
Let $\pi \in \mathfrak{D}^4_{2n}(1342)$. For all k, $1 \le k \le n$, if $\pi(2k) < 2k$ then $\pi(2k) = 2k-2$.
\end{lemma}

\begin{proof}
We will prove by contradiction that no deficiency at position $2k$ can have an entry less than or equal to $2k-4$.  So, let $\pi(2k) \le 2k-4$.  Since all odd entries are fixed points by Lemma \ref{lm:odd_fixed}, the two odd entries $2k-3$ and $2k-1$ are above and to the left of $\pi(2k)$, the triple $(2k-3, 2k-1,\pi(2k))$ at positions $(2k-3,2k-1,2k)$ would yield a 231-occurrence. 

By contradiction, the lemma is proved.
\end{proof}

Now we will prove Theorem \ref {1342_theorem}.

\begin{proof}[of Theorem \ref{1342_theorem}]
Given Lemma \ref{lm:odd_fixed}, all odd entries are fixed. Thus the even values are the only entries of interest, and moreover, all even values occur in even positions.  Furthermore, due to Lemma \ref{lm:def=2k-2}, a deficiency at position $2k$ must be the entry $2k-2$.

For $\pi\in\mathfrak{D}^4_{2n}(1342)=\mathfrak{D}^4_{2n}(231)$ consider a permutation $\pi_e\in S_n(231)$ such that $\pi_e(i)=j$ if $\pi(2i)=2j$ for $i,j\in[n]$. Then, for each $i\in[n]$, we have $\pi_e(i)\ge i-1$. Suppose that $\pi_e(1)=k$, then 231 avoidance of $\pi_e$ implies that the entries in $[1,k-1]$ occupy positions $[2,k]$. But $\pi_e(i)\ge i-1$ for all $i\in[n]$, so this can only occur if the values $1,\dots,k-1$ are consecutive and in increasing order. In other words, $\pi_e=(k,1,2,\dots,k-1,\tau+k)$, where $\tau\in S_{n-k}(231)$ satisfies the same conditions as $\pi_e$. 

Note also that the only nondeficiency among the first $k$ letters is the value $k$. Thus, we see inductively that the entire structure of $\pi_e$ is determined by the values of its nondeficiencies, one of which must be the entry $n$. Combining $\pi_e$ as above with the odd entries of $\pi$ (all fixed points), we see that $\pi$ also avoids $231$. Thus, $\pi_e$ (and hence, $\pi$) is determined by the choice of a subset $S \subseteq [1,n-1]$ such that if $k\in S$, then $2k$ is a nondeficiency of $\pi$. All of these choices are unrestricted, therefore the number of such choices is $2^{n-1}$. Moreover, the semilengths of the resulting blocks of the same form as $(1,2k,\pi')$ yield a composition of $n$ into nonzero parts. For example, $\pi = 16325478 \in \mathfrak{D}^4_{8}(1342)$ corresponds to the composition $4=3+1$ (see Figure \ref{fig:many8x8}). \qedhere


%\bigskip
%
%Now, let $\pi \in \mathfrak{D}^4_{2n}(1342)$ and consider $\pi(2)$.  If $\pi(2) = 2$, then $\pi = (1,2,\pi'+2)$ where $\pi' \in \mathfrak{D}^4_{2n-2}(1342)$.  So let $\pi(2) = 2k \ge 4$, for some $2 \le k \le n$.  Therefore all entries in $[2,2k-1]$ must precede all entries larger than $2k$, otherwise, we encounter a 231-occurrence. Then the block immediately to the right of $\pi(2)=2k$ consists of entries belonging to the interval $[2,2k-1]$ which are in positions $[3,2k]$.  Since the entries of interest are even values at even positions, we focus on even entries belonging to the set $[2,2k-2]$ which are in even positions in $[4,2k]$.
%
%Now, since $\pi(2) = 2k$ is an excedance, and since the entries in our block belong to the interval $[2,2k-2]$, the entry in position $2k$ must be a deficiency, namely the entry $2k-2$ by Lemma \ref{lm:def=2k-2}, i.e. $\pi(2k) = 2k-2$.  Therefore $\pi(2k-2)$ cannot be the entry $2k-2$, nor can it be an excedance since our entries belong to the interval $[2,2k-2]$.  Hence $\pi(2k-2)$ must be a deficiency, namely the entry $2k-4$,  i.e. $\pi(2k-2) = 2k-4$.  Continuing in this fashion, we have $\pi(2l) = 2l-2$ for $2 \le l \le k$ for the entire block with entries belonging to the set $[2,2k-2]$, which are in positions $[4,2k]$. More specifically, $\pi = (1,2k,\pi', \pi''+2k)$ where $\pi'=(3,2,5,4,7,6,\dots,2k-1,2k-2)$ is uniquely determined by the placement of $2k$, and $\pi'' \in \mathfrak{D}^4_{2n-2k}(1342)$. Notice also that if $k=1$, then $\pi'$ is empty. Moreover, the next even nondeficiency (a fixed point or an excedance) occurs at position $2k+2$. 
%
%Continuing recursively in this fashion, we see that the entire structure of $\pi$ is determined by the values of its even nondeficiencies.  Note that one of the nondeficiencies must be the entry $2n$, since $2n \ge \pi^{-1}(2n)$, the position of $2n$. Thus, $\pi$ is determined by the choice of a subset $S \subseteq [1,n-1]$ such that if $k\in S$, then $2k$ is a nondeficiency. All of these choices are unrestricted, therefore the number of such choices is $2^{n-1}$. Moreover, it is clear that the semilengths of the resulting blocks of the same form as $(1,2k,\pi')$ yield a composition of $n$ into nonzero parts.
%
%Pictorially, consider a $2n \times 2n$ board with dots in the $i$-th column from the left and $\pi(i)$-th row from the bottom.  We capture the dots in one or more $2k \times 2k$ squares $(1 \leq k \leq n)$ where each square starts with an odd left-to-right maximum (an entry $\pi(j)$ for which $\pi(j) > \pi(i)$ whenever $i < j$).  This coincides with the description of the block decomposition of $\pi$, since the second entry of each block (which is an even non-deficiency)) uniquely determines all entries of the block.
%
%Note that the diagonals of each $2k \times 2k$ square and the diagonal of the entire  $2n \times 2n$ board coincide, since all odd entries are fixed points.  Thus, if traveling from $(0,0)$ to $(2n,2n)$, the northeast corner of the last square must be $(2n,2n)$.  
%
%See Figure \ref{fig:big8x8_1342} corresponding to 
%For example, $\pi = 16325478 \in \mathfrak{D}^4_{8}(1342)$ corresponds to the composition $4=3+1$ (see Figure \ref{fig:many8x8}). \qedhere

%\vspace{0.25in}

%\begin{figure}[ht]
%\begin{center}
%\begin{tikzpicture}[scale=0.35] 
%%\draw[scale = 0.5, step = 2, very thin,color=gray] (0,0) grid (16,16);
%% \foreach \position in {(1,1), (5,5), (9,9), (13,13)} 
%%    \filldraw [scale = 0.5, color = black] \position circle (8pt);
%% \foreach \position in {(3,3), (7,11), (11,7), (15,15)} 
%%    \filldraw [scale = 0.5, color = red] \position circle (8pt);
%%\draw [line width = 3pt, color=blue] (0,0) rectangle (2,2);
%%\draw [line width = 3pt, color=blue] (2,2) rectangle (6,6);
%%\draw [line width = 3pt, color=blue] (6,6) rectangle (8,8);
%\draw[scale = 0.5, step = 2, very thin,color=gray] (0,0) grid (16,16);
% \foreach \position in {(1,1), (5,5), (9,9), (13,13)} 
%    \filldraw [scale = 0.5, color=black] \position circle (8pt);
% \foreach \position in {(3,11), (7,3), (11,7), (15,15)} 
%    \filldraw [scale = 0.5, color=red] \position circle (8pt);
%\draw [line width = 3pt, color=blue] (0,0) rectangle (6,6);
%\draw [line width = 3pt, color=blue] (6,6) rectangle (8,8);
%\end{tikzpicture} 
%\caption{Permutation diagram of $\pi = 16325478 \in \mathfrak{D}^4_{8}(1342)$.}
%\label{fig:big8x8_1342}
%\end{center}
%\end{figure}
%
%
%Note from Figure \ref{fig:big8x8_1342}, that at each odd left-to-right maximum (in this case $\pi(1)$ and $\pi(7)$) a new square begins, each odd entry is black to denote the fact that it must be fixed.  Also note that each even entry (which is colored red) may be fixed.  This process of capturing the dots from Dumont-4 permutations avoiding 1342 in the $2k \times 2k$ squares is the same as the various ways to split an integer into summands, also known as the number of compositions of the integer $n$ (since there are $n$ odd numbers in each permutation of length $2n$, and each square begins with an odd left-to-right maximum). And it is well known that the number of compositions of the integer $n$ is $2^{n-1}$.
\end{proof}

%This particular Dumont-4 permutation avoiding 1342 (see Figure \ref{fig:big8x8_1342}) would be the same as looking at the four odd entries in the squares and noting that this composition of the integer 4 would be $3+1$, as the first three odd entries occur in the bottom-left square, and the last odd entry occurs in the upper-right square.  For completion, see Example \ref {ex: D^4_8(1342)} below which outlines all $2^3 = 8$ Dumont-4 permutations of length 8 avoiding 1342, matching with each composition of the integer 4.

%\newpage

\begin{example} 
\label {ex: D^4_8(1342)}
Below are all eight Dumont-4 permutations of length $8$ avoiding $1342$, and the corresponding compositions of $4=8/2$.\\
\begin{center}

\begin{tikzpicture}[scale=0.4]

\begin{scope}
\draw[scale = 0.25, step = 2, very thin,color=gray] (0,0) grid (16,16);
 \foreach \position in {(1,1), (5,5), (9,9), (13,13)} 
    \filldraw [scale = 0.25, color=black] \position circle (8pt);
 \foreach \position in {(3,3), (7,7), (11,11), (15,15)} 
    \filldraw [scale = 0.25, color=red] \position circle (8pt);
\draw [line width = 1pt, color=blue] (0,0) rectangle (1,1);
\draw [line width = 1pt, color=blue] (1,1) rectangle (2,2);
\draw [line width = 1pt, color=blue] (2,2) rectangle (3,3);
\draw [line width = 1pt, color=blue] (3,3) rectangle (4,4);
\node [scale=0.6, below right] (0,0) {$4=1+1+1+1$};
\end{scope}

\begin{scope}[xshift=6cm]
\draw[scale = 0.25, step = 2, very thin,color=gray] (0,0) grid (16,16);
 \foreach \position in {(1,1), (5,5), (9,9), (13,13)} 
    \filldraw [scale = 0.25, color=black] \position circle (8pt);
 \foreach \position in {(3,3), (7,7), (11,15), (15,11)} 
    \filldraw [scale = 0.25, color=red] \position circle (8pt);
\draw [line width = 1pt, color=blue] (0,0) rectangle (1,1);
\draw [line width = 1pt, color=blue] (1,1) rectangle (2,2);
\draw [line width = 1pt, color=blue] (2,2) rectangle (4,4);
\node [scale=0.6, below right] (0,0) {$4=1+1+2$};
\end{scope}

\begin{scope}[xshift=12cm]
\draw[scale = 0.25, step = 2, very thin,color=gray] (0,0) grid (16,16);
 \foreach \position in {(1,1), (5,5), (9,9), (13,13)} 
    \filldraw [scale = 0.25, color=black] \position circle (8pt);
 \foreach \position in {(3,3), (7,11), (11,7), (15,15)} 
    \filldraw [scale = 0.25, color=red] \position circle (8pt);
\draw [line width = 1pt, color=blue] (0,0) rectangle (1,1);
\draw [line width = 1pt, color=blue] (1,1) rectangle (3,3);
\draw [line width = 1pt, color=blue] (3,3) rectangle (4,4);
\node [scale=0.6, below right] (0,0) {$4=1+2+1$};
\end{scope}

\begin{scope}[xshift=18cm]
\draw[scale = 0.25, step = 2, very thin,color=gray] (0,0) grid (16,16);
 \foreach \position in {(1,1), (5,5), (9,9), (13,13)} 
    \filldraw [scale = 0.25, color=black] \position circle (8pt);
 \foreach \position in {(3,3), (7,15), (11,7), (15,11)} 
    \filldraw [scale = 0.25, color=red] \position circle (8pt);
\draw [line width = 1pt, color=blue] (0,0) rectangle (1,1);
\draw [line width = 1pt, color=blue] (1,1) rectangle (4,4);
\node [scale=0.6, below right] (0,0) {$4=1+3$};
\end{scope}

\begin{scope}[xshift=0cm, yshift=-6cm]
\draw[scale = 0.25, step = 2, very thin,color=gray] (0,0) grid (16,16);
 \foreach \position in {(1,1), (5,5), (9,9), (13,13)} 
    \filldraw [scale = 0.25, color=black] \position circle (8pt);
 \foreach \position in {(3,7), (7,3), (11,11), (15,15)} 
    \filldraw [scale = 0.25, color=red] \position circle (8pt);
\draw [line width = 1pt, color=blue] (0,0) rectangle (2,2);
\draw [line width = 1pt, color=blue] (2,2) rectangle (3,3);
\draw [line width = 1pt, color=blue] (3,3) rectangle (4,4);
\node [scale=0.6, below right] (0,0) {$4=2+1+1$};
\end{scope}

\begin{scope}[xshift=6cm, yshift=-6cm]
\draw[scale = 0.25, step = 2, very thin,color=gray] (0,0) grid (16,16);
 \foreach \position in {(1,1), (5,5), (9,9), (13,13)} 
    \filldraw [scale = 0.25, color=black] \position circle (8pt);
 \foreach \position in {(3,7), (7,3), (11,15), (15,11)} 
    \filldraw [scale = 0.25, color=red] \position circle (8pt);
\draw [line width = 1pt, color=blue] (0,0) rectangle (2,2);
\draw [line width = 1pt, color=blue] (2,2) rectangle (4,4);
\node [scale=0.6, below right] (0,0) {$4=2+2$};
\end{scope}

\begin{scope}[xshift=12cm, yshift=-6cm]
\draw[scale = 0.25, step = 2, very thin,color=gray] (0,0) grid (16,16);
 \foreach \position in {(1,1), (5,5), (9,9), (13,13)} 
    \filldraw [scale = 0.25, color=black] \position circle (8pt);
 \foreach \position in {(3,11), (7,3), (11,7), (15,15)} 
    \filldraw [scale = 0.25, color=red] \position circle (8pt);
\draw [line width = 1pt, color=blue] (0,0) rectangle (3,3);
\draw [line width = 1pt, color=blue] (3,3) rectangle (4,4);
\node [scale=0.6, below right] (0,0) {$4=3+1$};
\end{scope}

\begin{scope}[xshift=18cm, yshift=-6cm]
\draw[scale = 0.25, step = 2, very thin,color=gray] (0,0) grid (16,16);
 \foreach \position in {(1,1), (5,5), (9,9), (13,13)} 
    \filldraw [scale = 0.25, color=black] \position circle (8pt);
 \foreach \position in {(3,15), (7,3), (11,7), (15,11)} 
    \filldraw [scale = 0.25, color=red] \position circle (8pt);
\draw [line width = 1pt, color=blue] (0,0) rectangle (4,4);
\node [scale=0.6, below right] (0,0) {$4=4$};
\end{scope}


\end{tikzpicture}

%$12345678 \leftrightarrow 1+1+1+1$ \hspace {1.3cm}  $12345876 \leftrightarrow 1+1+2$\\

%\vspace {1cm}
%
%\begin{tikzpicture}
%
%\end{tikzpicture}

%$12365478 \leftrightarrow 1+2+1$ \hspace {2cm}  $12385476 \leftrightarrow 1+3$\\


%\vspace {1cm}
%
%\begin{tikzpicture}
%
%\end{tikzpicture}

%$14325678 \leftrightarrow 2+1+1$ \hspace {2cm}  $14325876 \leftrightarrow 2+2$\\

%\vspace {1cm}
%
%\begin{tikzpicture}
%
%\end{tikzpicture}

%$16325478 \leftrightarrow 3+1$ \hspace {2cm}  $18325476 \leftrightarrow 4$\\

\begin{figure}[ht]
 \caption{Dumont-4 permutations of length $8$ avoiding $1342$.}
\label{fig:many8x8}
\end{figure}
\end{center}

\end{example}

%end of example

%\newpage

\subsubsection{$\mathfrak{D}^4_{2n}(1432)$ \label{321Dyck}}

As noted earlier, since the entry following $1$ is not $2$, it follows that $\mathfrak{D}^4_{2n}(1432)=\mathfrak{D}^4_{2n}(321)$.

\begin{theorem}
\label{thm:1432}
$|\mathfrak{D}^4_{2n}(1432)|=|\mathfrak{D}^4_{2n}(321)|=C_n$ for $n \ge 0$, where $C_n$ is the $n$-th Catalan number.
\label{321thm}
\end{theorem}

Since we are considering $321$-avoiding permutations, we will adapt a bijection of Krattenthaler \cite{Krattenthaler} from $S_n(321)$ to Dyck paths of semilength $n$.

\begin{proof}
Let $\pi \in \mathfrak{D}^4_{2n}(1432) = \mathfrak{D}^4_{2n}(321)$, and consider right-to-left minima (an entry $\pi(i)$ for which $\pi(i) < \pi(j)$ whenever $i < j$).  If an odd fixed point is a right-to-left minimum, then every entry to the left must have a smaller value, and every entry to the right must have a larger value, otherwise we encounter a 321-occurrence with the entry $2k+1$ serving as the ``2''.  In other words, if $\pi(2k+1) = 2k+1$, for some $0 \le k \le n-1$, then $\pi = (\pi', 2k+1, \pi''+2k+1)$ where $\pi' \in \mathfrak{D}^4_{2k}(321)$, and $(1,\pi''+1) \in \mathfrak{D}^4_{2n-2k}(321)$.  And, since the odd entry $\pi(2k+1) = 2k+1$ is a right-to-left minimum, the entry $2k+2$, which must be to the right of $2k+1$, is also a right-to-left minimum.

Consider a $2n \times 2n$ board with the dots in the $i$-th column from the left being in $\pi(i)$-th row from the bottom, for $1\le i\le 2n$. Now consider the dots which represent even right-to-left minima (solid red dots in Figure \ref{fig:big12x12_1432}), and lower them one cell down (hollow red dots in Figure \ref{fig:big12x12_1432}). Notice that, in particular, this will place a hollow red dot in each row with an odd right-to-left minimum.

Now travel along the cell boundaries from $(0, 0)$ to $(2n, 2n)$ in an \emph{East-North} fashion, using steps $(1,0)$ (\emph{east}) and $(0,1)$ (\emph{north}) and keeping all dots (both filled and hollow) to the left of the path, staying as close to the diagonal as possible. Equivalently, this is the path $P$ where the peaks (instances where an east step is followed directly by a north step) are exactly the bottom and right boundaries of the cells with hollow red dots.

For example, consider Figure \ref{fig:big12x12_1432} corresponding to $\pi = 1\ 3\ 5\ 2\ 6\ 4\ 7\ 8\ 9\ 11\ 12\ 10 \in \mathfrak{D}^4_{12}(321)$. Note that red dots are even right-to-left minima, the blue dots are odd right-to-left minima (which are fixed points), and the black dots are excedances.  

Suppose that $\pi$ has $k$ even non-excedances with (even) values $2=b_1<b_2<\dots<b_k$ at (even) positions $a_1<a_2<\dots<a_k=2n$, respectively. Also, let $b_{k+1}=2n+2$ and $a_0=0$. Then the \emph{runs} (maximal blocks) of east steps have lengths $a_i-a_{i-1}$ for $1\le i\le k$, and the runs of north steps have lengths $b_{i+1}-b_i$ for $1\le i\le k$.

Therefore, all runs of east and north steps in path $P$ are of even length (see Figure \ref{fig:big12x12_1432}, with $n=6$ and even non-excedance values of $2,4,8,10$ at positions $4,6,8,12$). Dividing the lengths of these runs in half, we obtain a Dyck path of semilength $n$ from $(0,0)$ to $(n,n)$.
% (see Figure \ref{fig:Dyck6x6_1432}).

%\vspace{0.25in}

\begin{figure}[!h]
\begin{center}
\begin{tikzpicture}[scale=0.175]
\draw[step = 2, very thin,color=gray] (0,0) grid (24,24);
\foreach \position in {(1,1), (3,5), (5,9), (9,11), (17,17), (19,21), (21, 23)} 
    \filldraw [color = black] \position circle (8pt);
\foreach \position in {(1,1), (13,13), (17,17)} 
    \filldraw [color = blue] \position circle (8pt);
\foreach \position in {(7,3), (11,7), (15,15), (23, 19)} 
    \filldraw [color = red] \position circle (8pt);
\foreach \position in {(7,1), (11,5), (15,13), (23, 17)} 
    \draw [color = red] \position circle (8pt);
\draw [->, line width = 1pt, color=red] (7,2.5) -- (7,1.5);
\draw [->, line width = 1pt, color=red] (11,6.5) -- (11,5.5);
\draw [->, line width = 1pt, color=red] (15,14.5) -- (15,13.5);
\draw [->, line width = 1pt, color=red] (23,18.5) -- (23,17.5);
\draw [line width = 3pt, color=black] (0,0) -- (8,0) -- (8,4) -- (12, 4) -- (12,12) -- (16,12) -- (16,16) -- (24,16) -- (24,24);
\draw [dashed, color = blue] (-0.5,-0.5) -- (24.5,24.5);
\end{tikzpicture} 
\caption{Permutation diagram of $\pi = 1\ 3\ 5\ 2\ 6\ 4\ 7\ 8\ 9\ 11\ 12\ 10 \in \mathfrak{D}^4_{12}(1432)$} 
\label{fig:big12x12_1432}
\end{center}
\end{figure}

%We can read each east step as a $\mathbf{u}$ and each north step as a $\mathbf{d}$.  Note, the diagonal of the permutation diagram in Figure \ref{fig:big12x12_1432} corresponds to the $x$-axis in the Dyck path in Figure \ref{fig:Dyck12x12_1432}, while $\mathbf{u}$ corresponds to step $(1,1)$ and $\mathbf{d}$ corresponds to step $(1,-1)$.
%
%\begin{figure}[!h]
%\begin{center}
%\begin{tikzpicture}[scale=0.4] 
%\foreach \position in {(0,0), (1,1), (2,2), (3,3), (4,4), (5,3), (6,2), (7,3), (8,4), (9,3), (10,2), (11,1), (12,0), (13,1), (14,2), (15,1), (16,0), (17,1), (18,2), (19,3), (20,4), (21,3), (22,2), (23,1), (24,0)} 
%    \filldraw [scale = 0.65, color = black] \position circle (6pt);
%%scale = 0.7071067
%\draw [scale = 0.65, line width = 2pt, color=black] (0,0) -- (4,4) -- (6,2) -- (8,4) -- (12,0) -- (14,2) -- (16,0) -- (20,4) -- (24,0);
%\draw [scale = 0.65, dashed] (-0.5,-0) -- (24.5,0);
%\end{tikzpicture} 
%\end{center}
%\caption{Dyck path corresponding to  $\pi = 1\ 3\ 5\ 2\ 6\ 4\ 7\ 8\ 9\ 11\ 12\ 10 \in \mathfrak{D}^4_{12}(1432)$}
%\label{fig:Dyck12x12_1432}
%\end{figure}
%
%%\vspace{0.25in}
%Note that the corresponding Dyck word is $\mathbf{uuuu\,dd\,uu\,dddd\,uu\,dd\,uuuu\,dddd}$. Dividing each run length in half, we obtain the string $\mathbf{uu\,d\,u\,dd\,u\,d\,uu\,dd}$ as in Figure \ref{fig:Dyck6x6_1432}.
%
%
%%Now to arrive at a Dyck path of semilength $n$, as desired, we map each Dyck path of semilength $2n$ to one and only one Dyck path of semilength $n$ by simply cutting each leg in half, where a leg may be defined as any number of steps between turns, or any number of consecutive steps in one direction.  This is attainable since every leg must be of even length.  We know this since in a Dumont-4 permutation, all even right-to-left minima are even entries at even positions.
%
%%Continuing with our running example, the Dyck path of semilength $2n$ corresponding to $\pi = 1\ 3\ 5\ 2\ 6\ 4\ 7\ 8\ 9\ 11\ 12\ 10 \in \mathfrak{D}^4_{12}(1432)$ maps to the following Dyck path of semilength $n$.  See Figure \ref{fig:Dyck6x6_1432}.
%
%%\newpage
%\vspace{0.25in}
%
%\begin{figure}[!h]
%\begin{center}
%\begin{tikzpicture}[scale=0.4] 
%
%\begin{scope}
%
%%scale = 0.7071067
%\draw [scale = 0.65, line width = 2pt, color=black] (0,0) -- (4,4) -- (6,2) -- (8,4) -- (12,0) -- (14,2) -- (16,0) -- (20,4) -- (24,0);
%\draw [scale = 0.65, dashed] (-0.5,-0) -- (24.5,0);
%\foreach \position in {(1,1), (3,3), (5,3), (7,3), (9,3), (11,1), (13,1), (15,1), (17,1), (19,3),(21,3), (23,1)} 
%    \filldraw [scale = 0.65, color = black] \position circle (6pt);
%\foreach \position in {(0,0), (2,2), (4,4), (6,2), (8,4), (10,2), (12,0), (14,2), (16,0), (18,2), (20,4), (22,2), (24,0)} 
%    \draw [scale = 0.65, color = black, fill=white] \position circle (5pt);
%\end{scope}
%
%\begin{scope}[xshift=18cm,scale=2]
%%scale = 0.7071067
%\draw [scale = 0.65, line width = 2pt, color=black] (0,0) -- (2,2) -- (3,1) -- (4,2) -- (6,0) -- (7,1) -- (8,0) -- (10,2) -- (12,0);
%\draw [scale = 0.65, dashed] (-0.5,-0) -- (12.5,0);
%\foreach \position in {(0,0), (1,1), (2,2), (3,1), (4,2), (5,1), (6,0), (7,1), (8,0), (9,1), (10,2), (11,1), (12,0)} 
%   \draw [scale = 0.65, color = black, fill=white] \position circle (2.5pt);
%\end{scope}
%
%\end{tikzpicture} 
%
%%\vspace{0.5in}
%%
%%\begin{tikzpicture}[scale=0.5] 
%%
%%
%%\end{tikzpicture} 
%
%\caption{Dyck path of semilength $2n$ mapping to Dyck path of semilength $n$}
%\label{fig:Dyck6x6_1432}
%\end{center}
%\end{figure}

Therefore, the number of Dumont-4 permutations of length $2n$ avoiding $321$ is the same as the number of Dyck paths of semilength $n$, i.e. the $n$th Catalan number, $C_n$.  This ends the proof.
\end{proof}


%Note, the corresponding Dyck word would reflect each consecutive run of \textbf{u}'s and \textbf{d}'s cut in half as well; our Dyck word is now $\mathbf{uu\,d\,u\,dd\,u\,d\,uu\,dd}$.  We are now at our desired length of a Dyck path travelling from $(0,0)$ to $(n,n)$ in only $(1,1)$ or $(-1,1)$ steps, and not going below the $x$-axis.  This sequence, for all $n \ge 0$, enumerate the well known Catalan numbers. 


%\vspace{0.25in}

%%% Section - Enumerating Dumont-4 permutations avoiding certain permutations of length four %%%

\subsection {Enumerating Dumont-4 permutations avoiding certain permutations of length four}

Now that all three Dumont-4 permutations of length four have been avoided by Dumont-4 permutations of length $2n$, we will look at three other permutations of length four starting with $1$, namely $1324$, $1243$, and $1423$.  Also note, excluding the entry ``$1$", the three permutations in the previous section coupled with the three permutations in this section constitute all of $S_3$.

\subsubsection{$\mathfrak{D}^4_{2n}(1324)$}

As noted earlier, since the entry following $1$ is not $2$, it follows that $\mathfrak{D}^4_{2n}(1324)=\mathfrak{D}^4_{2n}(213)$.

\begin{theorem}
$|\mathfrak{D}^4_{2n}(1324)| = 2 \binom{n}{2} + 1 = n^2 - n + 1$, \text{for} $n \ge 0$.
\end{theorem}

\begin{proof}
Consider the rightmost value of $\pi\in\mathfrak{D}^4_{2n}(1324)$. If the entry $2n$ occurs in the last position, then $\pi = 1,2,3,4,\dots,2n-1,2n$.  This permutation accounts for the ``$1$" in $2\binom{n}{2} + 1$.

If the last entry is not $2n$, then it must be a deficiency, so by definition of a Dumont-4 permutation, it must be even, say $2k$, where $1 \le k \le n-1$. 

Note that $\pi(2n-1)=2n-1$ or $\pi(2n-1)=2n$, and therefore all entries below and to the left of $\pi(2n-1)$ must be in increasing order to avoid pattern $213$. Therefore, all entries in $[1,2n-1]\setminus\{2k\}$, i.e. all entries of $\pi$ except $2n$ and $2k$ are in increasing order.

Moreover, the value $2k-1$ can only occur in position $2k-1$ or smaller, so for the values $1,2,\dots,2k-1$ to form an increasing subsequence of $\pi$, all of them must be fixed points. Thus, the value $2n$ must occur in some position $l\in[2k,2n-1]$. In fact, each choice of $\pi(2n)=2k$ and $\pi^{-1}(2n)=l\in[2k,2n-1]$ yields a unique permutation $\pi\in\mathfrak{D}^4_{2n}(1324)$.

Thus, 
%using the block decomposition of Mansour and Vainshtein \cite{MV_Restricted}, 
we can write $\pi=(1,\pi',\rho,2n,\sigma,2k)$, where $\pi'= (1, 2, \dots, 2k-1)$, $\rho=(2k+1, 2k+2, \dots, l)$, and $\sigma=(l+1,l+2, \dots, 2n-1)$.


%
%If the last entry is not $2n$, then it must be a deficiency, so by definition of a Dumont-4 permutation, it must be even, say $2k$, where $1 \le k \le n-1$.  Using the block decomposition method of Mansour and Vainshtein in \cite{MV_Restricted}, we have $\pi = (1,\pi', \pi'', 2k)$, where 
%%and referenced and defined in \cite{MV_Chebyshev} by them as well
%\begin{itemize}
%\item $\pi'$ is a permutation of the set $[2, 2k-1]$,
%\item $\pi''$ is a permutation of the set $[2k+1,2n]$, and
%\item both $\pi'$ and $\pi''$ avoid 213.
%\end{itemize}
%
%%\vspace{0.5in}
%
%\begin{center}
%\begin{tikzpicture}[scale=0.5]
% \foreach \position in {(7.5,2.5), (-0.5,-0.5)}
%    \filldraw [color = red] \position circle (4pt);
% \node [right] at (7.7,2.5) {$2k$};
% \node [left] at (-0.7,-0.5) {$1$};
%\draw [line width = 2pt, color = red] (0,0) rectangle (2,2);
% \node at (1,1) 
%	{$\begin{aligned} 
%			\pi'
%     		\end{aligned}$};
%\draw [line width = 2pt, color = red] (2,3) rectangle (7,8);
% \node at (4.5,5.5) 
%	{$\begin{aligned} 
%			\pi''
%     		\end{aligned}$};
%\end{tikzpicture} 
%\begin{figure}[ht]
%\caption{Block decomposition for $\pi \in \mathfrak{D}^4_{2n}(1324)$ when $2k$ is in position $2n$.}
%\label{fig:block1324_0}
%\end{figure}
%\end{center} 
%
%\noindent Note that the block representing $\pi''$ must contain $2n$.  Now, reconsidering the entry $2n$, again, if it is the last entry of the block representing $\pi''$ the resulting permutation would be increasing on the diagonal of that block.  The block representing $\pi'$ would be increasing as well.  Thus we have, $$\pi = (\pi', \pi'', 2k) = \{1, 2, 3, \dots, 2k-1, 2k+1, 2k+2, \dots, 2n-1, 2n, 2k\}.$$  
%
%Now, since $2n$ is an excedance, it could be in any position in the block representing $\pi''$, say position $l\in[2k+1,2n-1]$. 
%
%Note that $\pi(2n-1)=2n-1$ or $2n$ and that all entries below and to the left of $\pi(2n-1)$ must be in increasing order to avoid 213. Therefore, all entries in $[1,2n-1]\setminus\{2k\}$, i.e. all entries of $\pi$ except $2n$ and $2k$, are in increasing order, which determines their positions uniquely.
%
%This forces another decomposition inside the block representing $\pi''$.  In other words, we have $\pi'' = (\rho, 2n, \sigma)$ where,
%
%\begin{itemize}
%\item $\pi'= 1, 2, \dots, 2k-1$,
%\item $\rho=2k+1, 2k+2, \dots, l$, and
%\item $\sigma=l+1,l+2, \dots, 2n-1$.
%\end{itemize}

%\vspace{0.25in}

%\begin{figure}[!h]
%\begin{center}
%\begin{tikzpicture}[scale=0.65]
% \foreach \position in {(7.5,2.5), (-0.5,-0.5)}
%    \filldraw [color = red] \position circle (4pt);
%
%\draw [line width = 2pt, color = red] (0,0) rectangle (2,2);
% \node at (1,1) {$\pi'$};
%
%\draw [line width = 2pt, color = red] (2,3) rectangle (7,8);
% \node at (3,7) {$\pi''$};
%
% \foreach \position in {(4.5,7.5)}
%    \filldraw [color = blue] \position circle (4pt);
%
%\draw [line width = 2pt, color = blue, dashed] (2.1,3.1) rectangle (4,5);
% \node at (3,4) {$\rho$};
%
%\draw [line width = 2pt, color = blue, dashed] (5,5) rectangle (6.9,7);
%\node at (6,6) {$\sigma$};
%
%\node at (5,7.5) {$2n$};
%\node at (8,2.5) {$2k$};
%\node [left] at (-0.7,-0.5) {$1$};
%\end{tikzpicture} 
%\caption{Block decomposition for $\pi \in \mathfrak{D}^4_{2n}(1324)$.
%% when $\pi(2n)=2k$ and $\pi^{-1}(2n)=l\in[2k,2n-1]$
%}
%\label{fig:block1324_0_2}
%\end{center} 
%\end{figure}

%Again, all other entries within their blocks must be increasing on their respective diagonals.  If not, any other configuration would result in a 213-occurrence.  Thus, all Dumont-4 permutations avoiding 213 can be determined by the occurrence of the entry $2k$ in the last position, and the occurrence of the entry $2n$ in the $l$-th position, which again must be in the block representing $\pi''$.  So if we look at the last position, starting at the top (the entry $2n$), and move it down to an allowable entry (any even entry less than $2n$), then $2k$ can take the value of $2n-2, 2n-4, \dots, 4,$ and $2$.  Each movement downward, after the top-right location, creates two more choices for the $l$-th position of the entry $2n$ in the block represented by $\pi''$.

When $\pi(2n)=2k<2n$, the number of possible choices for the position $l$ of $2n$ is the number of elements in $[2k,2n-1]$, i.e. $2n-2k$. Since $1\le k\le n-1$, the total number of Dumont-4 permutations avoiding 1324 is
\[
|\mathfrak{D}^4_{2n}(1324)| = 1 + \sum_{k=1}^{n-1}(2n-2k)
= 1 + 2\sum_{k=1}^{n-1}(n-k)
= 1 + 2\sum_{k=1}^{n-1}k
= 1+ 2 \binom {n}{2}
= n^2-n+1. \qedhere
\]
\end{proof}


\begin{example}
\label{example:D^4_16(1324)}
A Dumont-4 permutation of length $16$ avoiding $1324$ can be obtained by placing the entry $6$ in position $16$, and the entry $16$ in position $10$. All other entries must be in increasing order.  
%We leave out the first entry ``1" from the block in preparation for the next section.  
See Figure \ref{fig:big16x16} for an example.

\begin{figure}[!ht]
\begin{center}
\begin{tikzpicture}[scale = 0.4]
\draw[very thin, color = gray, scale = 0.7] (0,0) grid (16,16);
 \foreach \position in {(1,1), (3,3), (5,5), (7,7), (9,9), (11,13), (13,15), (15,17), (17,19), (21,21), (23,23), (25, 25), (27, 27), (29, 29)}
    \filldraw [scale = 0.35, color = black] \position circle (8pt);
 \foreach \position in {(31,11)} 
    \filldraw [scale = 0.35, color = red] \position circle (8pt);
 \foreach \position in {(19,31)} 
    \filldraw [scale = 0.35, color = blue] \position circle (8pt);
\draw [line width = 2pt, color = red, scale = 0.7] (1,1) rectangle (5,5);
\draw [line width = 2pt, color = red, scale = 0.7] (5,6) rectangle (15,16);
\draw [line width = 2pt, color = blue, dashed, scale = 0.7] (5.1,6.1) rectangle (9,10);
\draw [line width = 2pt, color = blue, dashed, scale = 0.7] (10, 10) rectangle (14.9,15);
\end{tikzpicture} 
	\caption {Permutation diagram of
	$\pi = 	1\ 2\ 3\ 4\ 5\ 7\ 8\ 9\ 10\ 16\ 11\ 12\ 13\ 14\ 15\ 6 
			\in \mathfrak{D}^4_{16}(1324)$}
	\label{fig:big16x16}
\end{center}
\end{figure}

\end{example}

\subsubsection{$\mathfrak{D}^4_{2n}(1243)$}

\begin{theorem} \label{thm:1243}
$|\mathfrak{D}^4_{2n}(1243)| = |\mathfrak{D}^4_{2n}(1324)| = 2 \binom{n}{2} + 1 = n^2 - n + 1$, \text{for} $n \ge 0$.
\end{theorem}

In other words, patterns $1324$ and $1243$ are Wilf-equivalent on Dumont permutations of the fourth kind.

\begin{proof}
This enumeration is the same as that of $\mathfrak{D}^4_{2n}(1324)$ since the removal of the entry $1$ from patterns $1324$ and $1243$ yields the patterns $213$ and $132$ that are each other's reflections with respect to the antidiagonal.  That is, $213$ and $132$ are inverses of reversals of complements of each other, and thus their respective avoidance classes are enumerated by the same sequence.

To obtain a permutation in $\mathfrak{D}^4_{2n}(1243)$ from a permutation in $\mathfrak{D}^4_{2n}(1324)$, remove the entry ``$1$" and reflect the remaining entries from  about the antidiagonal. This maps blocks onto blocks, with their diagonals mapping onto the diagonals of the images of those blocks.  Lastly, add $1$ to every entry and prepend the value $1$ that was removed. See Figure \ref{fig:big16x16_irc} for an example. Note that reflection about the antidiagonal applied above implies that any permutation $\pi\in\mathfrak{D}^4_{2n}(1243)$ is uniquely given by $\pi^{-1}(2)$ and $\pi(2)\in[2,\pi^{-1}(2)]$.
\end{proof}

\begin{example}
\label{example:D^4_16(1243)}
Reflecting the 1324-avoiding permutation in Figure \ref{fig:big16x16} as in Theorem \ref{thm:1243}, we obtain a 1243-avoiding Dumont-4 permutation.
%
%So, after removing the ``1", and subtracting 1 from each entry remaining, instead of placing the entry 5 in the last position (fifth from the bottom), we place the entry 1 in the bottom row (fifth from the right).  Also, instead of the entry 15 occuring in the top row (seventh from the right), we place the entry 7 in the first position (seventh from the bottom).
%
\begin{figure}[!h]
\begin{center}
\begin{tikzpicture}[scale = 0.4]
\draw[very thin, color = gray, scale = 0.7] (0,0) grid (16,16);
 \foreach \position in {(1,1), (5,5), (7,7), (9,9), (11,11), (13,13), (15,17), (17,19), (19,21), (21,23), (25, 25), (27, 27), (29, 29), (31,31)}
    \filldraw [scale = 0.35, color = black] \position circle (8pt);
 \foreach \position in {(23,3)} 
    \filldraw [scale = 0.35, color = red] \position circle (8pt);
 \foreach \position in {(3,15)} 
    \filldraw [scale = 0.35, color = blue] \position circle (8pt);
\draw [line width = 2pt, color = red, scale = 0.7] (12,12) rectangle (16,16);
\draw [line width = 2pt, color = red, scale = 0.7] (1,2) rectangle (11,12);
\draw [line width = 2pt, color = blue, dashed, scale = 0.7] (2.1,2.1) rectangle (7,7);
\draw [line width = 2pt, color = blue, dashed, scale = 0.7] (7, 8) rectangle (10.9,11.9);
\end{tikzpicture} 
	\caption {Permutation diagram of
	$\rho = 	1\ 8\ 3\ 4\ 5\ 6\ 7\ 9\ 10\ 11\ 12\ 2\ 13\ 14\ 15\ 16 		
					\in \mathfrak{D}^4_{16}(1243)$.}
	\label{fig:big16x16_irc}
\end{center}
\end{figure}
\end{example}

\subsubsection{$\mathfrak{D}^4_{2n}(1423)$}

As noted earlier, since the entry following $1$ is not $2$, it follows that $\mathfrak{D}^4_{2n}(1423)=\mathfrak{D}^4_{2n}(312)$.

Since our result for this pattern involves a continued fraction, we need a remark on notation.

\begin{notation}
For convenience, the following notation will be used for continued fractions:
\[
\cfrac{a}{\alpha \pm \cfrac{b}{\beta \pm \cfrac{c}{\gamma \pm \ddots}}} \quad = \quad 
\frac{a}{\alpha} \,\, \genfrac{}{}{0pt}{}{}{\pm} \,\, \frac{b}{\beta} \,\, \genfrac{}{}{0pt}{}{}{\pm} \,\, \frac{c}{\gamma} \,\, \genfrac{}{}{0pt}{}{}{\pm} \,\, \cdots
\]
\end{notation}

%\vspace{0.25in}

Additionally, we will need to use the truncations of the even and odd parts of the Catalan generating function $C(z)=\sum_{n=0}^{\infty}{C_n z^n}$. Define 
\[
C_{e,2m}=C_{e,2m}(z) =  \sum_{i=0}^{m} C_{2i}z^{2i} 
%= \frac{C(z) + C(-z)}{2}
\quad \text{and} \quad 
C_{o,2m+1}=C_{o,2m+1}(z) =  \sum_{i=0}^{m} C_{2i+1}z^{2i+1},
%= \frac{C(z) - C(-z)}{2}
\]
where we set these functions equal to $0$ when $m<0$.

\begin{theorem}
For all $k \ge 1$, the ordinary generating function for $|\mathfrak{D}^4_{2n}(1423)|$ is $(R_1/z)\circ\sqrt{z}=(zR_1/z^2)\circ\sqrt{z}$, where the sequence of functions $zR_{2k+1}$, $k\ge 0$, satisfies the recurrence relation
\begin{equation} \label{eq:d4-312-ogf}
zR_{2k+1} = \frac{z^2C_{e,2k}}{(1-zC_{o,2k-1})^2} \,\,
 \genfrac{}{}{0pt}{}{}{-} \,\,
\frac{z^2C_{e,2k}}{1-zC_{o,2k+1}} \,\,
 \genfrac{}{}{0pt}{}{}{-} \,\,
\frac{z^2C_{e,2k}^2}{1} \,\,
 \genfrac{}{}{0pt}{}{}{-} \, \,
zR_{2k+3}.
\end{equation} 
\end{theorem}

This yields the generating function for $|\mathfrak{D}^4_{2n}(1423)|$ in the form of a continued fraction.

%\vspace{0.25in}

We consider a Dumont-4 permutation avoiding $1423$, and we are interested in the parity of the smallest entry in each block using block decomposition after the ``$1$''.  Note, since every Dumont-4 permutation begins with the entry ``$1$'', avoiding $1423$ is the same as avoiding $312$ by the Dumont-4 permutation with the ``$1$'' removed from the beginning of the permutation.  Therefore we analyze the block decomposition of the resulting permutation diagram, ignoring the initial block of ``$1$''.

To analyze the blocks resulting from the iterations of the block decompositon, we will need an auxiliary parameter, namely, the length $m$ of the maximal contiguous segment of allowed cells in the bottom row. Our cases are further subdivided according to the parity of $m$ and $n-m$, where $n$ is the dimension of the (square) block. We will refer to the cells on the line $y=x-m$ as the \emph{$m$-th subdiagonal} of a board and call the cells below the $m$-th subdiagonal \emph{$m$-deficiencies} (so, deficiencies as defined earlier are 0-deficiencies in this terminology).

Now, if the bottom row of a block is an odd row in the starting diagram of a Dumont-4 permutation then there is no $m$-deficiency in that bottom row.  If the bottom row of a block is an even row in the starting diagram, then there may be an $m$-deficiency in that bottom row.  Given the definition of the Dumont-4 permutations, in the boards resulting from the repeated block decomposition, all cells  above or to the left of the $m$-th subdiagonal are allowed, whereas the positions and values of the possible $m$-deficiencies are parity-restricted as described below:
\begin{description}
\item[EE blocks:] $m=2k$, $n-m$ is even (so $n$ is even), and all $m$-deficiencies are even values in odd positions;
\item[NE blocks:] $m=2k+1$, $n-m$ is even (so $n$ is odd), and all $m$-deficiencies are even values in even positions;
\item[EN blocks:] $m=2k$, $n-m$ is odd (so $n$ is odd), and all $m$-deficiencies are odd values in even positions;
\item[NN blocks:] $m=2k+1$, $n-m$ is odd (so $n$ is even), and all $m$-deficiencies are odd values in odd positions.
\end{description}

We denote the blocks mnemonically using ``E'' for even and ``N'' for ``not even,'' i.e. odd. However, for convenience in working with generating functions, we also let the $EE$ (resp. $NE, EN$, and $NN$) block with parity-restricted positions and values of $m$-deficiencies be represented by the generating function $P_m=P_m(z)$ (resp. $R_m=R_m(z),S_m=S_m(z)$, and $T_m=T_m(z)$), and refer to that block as a $P$-board (resp. $R$-board, $S$-board, and $T$-board).  See Figure \ref{312_odd_even row}, where blue dots are odd positions and red dots are even positions in the starting $\mathfrak{D}^4$-permutation diagram.

%See Figure \ref{312_odd row} and Figure \ref{312_even row}, where blue dots are odd positions and red dots are even positions.

%\vspace{0.25in}
%\newpage

\begin{figure}[!ht]
\begin{center}
\begin{tikzpicture}[scale = 0.45]
\newcommand{\mysize}{0.5}

\begin{scope}
\draw[very thin, scale = \mysize, color = gray] (0,0) grid (14,14);
\foreach \x in {6} \foreach \y in {0}
\draw [line width = 2pt, scale = \mysize, color = black] (0,0) -- (\x,\y) -- (\x,\y+1) -- (\x+1,\y+1) -- (\x+1,\y+2) -- (\x+2,\y+2) -- (\x+2,\y+3) -- (\x+3,\y+3) -- (\x+3,\y+4) -- (\x+4,\y+4) -- (\x+4,\y+5) -- (\x+5,\y+5) -- (\x+5,\y+6) -- (\x+6,\y+6) -- (\x+6,\y+7) -- (\x+7,\y+7) -- (\x+7,\y+8) -- (\x+8,\y+8) -- (\x+8,\y+14) -- (0,\y+14) -- cycle;
\foreach \x in {8,10,12} \foreach \y in {1}
{\draw [line width = 2pt, scale = 0.5, color = black] (\x,\y) rectangle (\x+1, \y+1);}
\foreach \x in {10,12} \foreach \y in {3}
{\draw [line width = 2pt, scale = 0.5, color = black] (\x,\y) rectangle (\x+1, \y+1);}
\foreach \x in {12} \foreach \y in {5}
{\draw [line width = 2pt, scale = 0.5, color = black] (\x,\y) rectangle (\x+1, \y+1);}
\foreach \position in {(0.5,0.5), (2.5, 0.5), (4.5,0.5)} 
    \filldraw [scale = \mysize, color = red] \position circle (4pt);
 \foreach \position in {(1.5,0.5), (3.5, 0.5), (5.5,0.5)} 
    \filldraw [scale = \mysize, color = blue] \position circle (4pt);
\node[scale=0.8] at (1.5,-0.5) {$2k$};
\node[scale=0.8, rotate=90] at (7.5, 5.5) {$2k$};
\node[scale=0.8] at (3.5, 8.5) {\underline{EE block ($P$-board)}};
\end{scope}

\begin{scope}[xshift=9cm]
\draw[very thin, scale = \mysize, color = gray] (0,0) grid (15,15);
\foreach \x in {7} \foreach \y in {0}
\draw [line width = 2pt, scale = \mysize, color = black] (0,0) -- (\x,\y) -- (\x,\y+1) -- (\x+1,\y+1) -- (\x+1,\y+2) -- (\x+2,\y+2) -- (\x+2,\y+3) -- (\x+3,\y+3) -- (\x+3,\y+4) -- (\x+4,\y+4) -- (\x+4,\y+5) -- (\x+5,\y+5) -- (\x+5,\y+6) -- (\x+6,\y+6) -- (\x+6,\y+7) -- (\x+7,\y+7) -- (\x+7,\y+8) -- (\x+8,\y+8) -- (\x+8,\y+15) -- (0,\y+15) -- cycle;
\foreach \x in {9,11,13} \foreach \y in {1}
{\draw [line width = 2pt, scale = \mysize, color = black] (\x,\y) rectangle (\x+1, \y+1);}
\foreach \x in {11,13} \foreach \y in {3}
{\draw [line width = 2pt, scale = \mysize, color = black] (\x,\y) rectangle (\x+1, \y+1);}
\foreach \x in {13} \foreach \y in {5}
{\draw [line width = 2pt, scale = \mysize, color = black] (\x,\y) rectangle (\x+1, \y+1);}
\foreach \position in {(0.5,0.5), (2.5, 0.5), (4.5,0.5), (6.5,0.5)} 
    \filldraw [scale = \mysize, color = blue] \position circle (4pt);
 \foreach \position in {(1.5,0.5), (3.5, 0.5), (5.5,0.5)} 
    \filldraw [scale = \mysize, color = red] \position circle (4pt);
\node[scale=0.8] at (1.75,-0.5) {$2k+1$};
\node[scale=0.8, rotate=90] at (8, 5.75) {$2k+1$};
\node[scale=0.8] at (3.6, 8.5) {\underline{NE block ($R$-board)}};
\end{scope}

%\end{tikzpicture} 
%\caption{Blocks whose bottom row was odd in the starting $\mathfrak{D}^4$-permutation diagram}
%\label{312_odd row}
%\end{figure}
%
%%\newpage
%%\vspace{0.5in}
%
%\begin{figure}[!h]
%\begin{tikzpicture}[scale = 0.5]
%\newcommand{\mysize}{0.5}

\begin{scope}[xshift=18cm]
\draw[very thin, scale = \mysize, color = gray] (0,0) grid (15,15);
\foreach \x in {6} \foreach \y in {0}
\draw [line width = 2pt, scale = \mysize, color = black] (0,0) -- (\x,\y) -- (\x,\y+1) -- (\x+1,\y+1) -- (\x+1,\y+2) -- (\x+2,\y+2) -- (\x+2,\y+3) -- (\x+3,\y+3) -- (\x+3,\y+4) -- (\x+4,\y+4) -- (\x+4,\y+5) -- (\x+5,\y+5) -- (\x+5,\y+6) -- (\x+6,\y+6) -- (\x+6,\y+7) -- (\x+7,\y+7) -- (\x+7,\y+8) -- (\x+8,\y+8) -- (\x+8,\y+9) -- (\x+9,\y+9) -- (\x+9,\y+15) -- (0,\y+15) -- cycle;
\foreach \x in {7,9,11,13} \foreach \y in {0}
{\draw [line width = 2pt, scale = \mysize, color = black] (\x,\y) rectangle (\x+1, \y+1);}
\foreach \x in {9,11,13} \foreach \y in {2}
{\draw [line width = 2pt, scale = \mysize, color = black] (\x,\y) rectangle (\x+1, \y+1);}
\foreach \x in {11, 13} \foreach \y in {4}
{\draw [line width = 2pt, scale = \mysize, color = black] (\x,\y) rectangle (\x+1, \y+1);}
\foreach \x in {13} \foreach \y in {6}
{\draw [line width = 2pt, scale = \mysize, color = black] (\x,\y) rectangle (\x+1, \y+1);}
\foreach \position in {(0.5,0.5), (2.5, 0.5), (4.5,0.5)} 
    \filldraw [scale = \mysize, color = blue] \position circle (4pt);
 \foreach \position in {(1.5,0.5), (3.5, 0.5), (5.5,0.5), (7.5, 0.5), (9.5, 0.5), (11.5, 0.5), (13.5,0.5)} 
    \filldraw [scale = \mysize, color = red] \position circle (4pt);
\node[scale=0.8] at (1.5,-0.5) {$2k$};
\node[scale=0.8, rotate = 90] at (8, 6) {$2k$};
\node[scale=0.8] at (3.5, 8.5) {\underline{EN block ($S$-board)}};
\end{scope}

\begin{scope}[xshift=27cm]
\draw[very thin, scale = \mysize, color = gray] (0,0) grid (14,14);
\foreach \x in {5} \foreach \y in {0}
\draw [line width = 2pt, scale = \mysize, color = black] (0,0) -- (\x,\y) -- (\x,\y+1) -- (\x+1,\y+1) -- (\x+1,\y+2) -- (\x+2,\y+2) -- (\x+2,\y+3) -- (\x+3,\y+3) -- (\x+3,\y+4) -- (\x+4,\y+4) -- (\x+4,\y+5) -- (\x+5,\y+5) -- (\x+5,\y+6) -- (\x+6,\y+6) -- (\x+6,\y+7) -- (\x+7,\y+7) -- (\x+7,\y+8) -- (\x+8,\y+8) -- (\x+8,\y+9) -- (\x+9,\y+9) -- (\x+9,\y+14) -- (0,\y+14) -- cycle;
\foreach \x in {6,8,10,12} \foreach \y in {0}
{\draw [line width = 2pt, scale = \mysize, color = black] (\x,\y) rectangle (\x+1, \y+1);}
\foreach \x in {8,10,12} \foreach \y in {2}
{\draw [line width = 2pt, scale = \mysize, color = black] (\x,\y) rectangle (\x+1, \y+1);}
\foreach \x in {10,12} \foreach \y in {4}
{\draw [line width = 2pt, scale = \mysize, color = black] (\x,\y) rectangle (\x+1, \y+1);}
\foreach \x in {12} \foreach \y in {6}
{\draw [line width = 2pt, scale = \mysize, color = black] (\x,\y) rectangle (\x+1, \y+1);}
\foreach \position in {(0.5,0.5), (2.5, 0.5), (4.5,0.5), (6.5,0.5), (8.5, 0.5), (10.5,0.5), (12.5, 0.5)} 
    \filldraw [scale = \mysize, color = red] \position circle (4pt);
 \foreach \position in {(1.5,0.5), (3.5,0.5)} 
    \filldraw [scale = \mysize, color = blue] \position circle (4pt);
\node[scale=0.8] at (1.25,-0.5) {$2k+1$};
\node[scale=0.8, rotate=90] at (7.5, 5.75) {$2k+1$};
\node[scale=0.8] at (3.6, 8.5) {\underline{NN block ($T$-board)}};
\end{scope}

\end{tikzpicture} 
\end{center}

\caption{Blocks whose bottom row was odd (even) in the starting $\mathfrak{D}^4$-permutation diagram}
\label{312_odd_even row}
\end{figure}


\begin{proof}
We now produce the generating function for the enumeration sequence of $\mathfrak{D}^4_{2n} (1423) = \mathfrak{D}^4_{2n} (312)$.  To begin, note that a $P$-board and $T$-board may be empty, while an $R$-board and $S$-board cannot be empty.  This is due to the fact that $P$-boards and $T$-boards have even dimensions whereas  $R$-boards and $S$-boards have odd dimensions. 

Consider the recurrence relations for the generating functions corresponding to each of the four blocks. We use the position of the entry $1$ in the block to produce those. Since a permutation in each block avoids $312$, any value to the left of $1$ must be less that any value to the right of $1$. For the $P$-board with $m=2k$, $k\ge 1$, we have the generating function $P_{2k}(z)$, the recurrence relation for which has three summands:
\begin{itemize}
\item The $P$-board may be empty, which corresponds to the summand ``$1$".  
\item If the $P$-board is nonempty, consider the entry in the bottom row. If it is in an odd position $\le 2k-1$, the factor of $z$ corresponds to the bottom row entry, in this case a blue dot on the $P$-board in Figure~\ref{312_odd_even row}. The block to the left of the blue dot must be a square of even dimension $\le 2k-2$ with all cells allowed, and the block to the right of the blue dot is an $S$-board with $m=2k$, correponding to the generating function $S_{2k}$.  This yields the summand ${\color{red}z C_{e,2k-2} S_{2k}}$.
\item Lastly, consider the entry in the bottom row that is in an even position $\le 2k$.  Again, the factor $z$ corresponds to the bottom row entry, in this case is a red dot on the $P$-board in Figure~\ref{312_odd_even row}. The block to the left of the red dot must be a square of odd dimension $\le 2k-1$ with all cells allowed, and the block to the right of the red dot is a $P$-board with $m=2k$, corresponding to the generating function $P_{2k}$.  This yields the summand ${\color{blue}z C_{o,2k-1}P_{2k}}$.
\end{itemize}

%\noindent
Thus, the recurrence formula for our generating function $P_{2k}$ is given by 
\[
P_{2k} = 1 + {\color{red}z C_{e,2k-2} S_{2k}} + {\color{blue}z C_{o,2k-1}P_{2k}}.
\]

%\newpage

%\noindent 
We find the remaining recurrence formulas for $R_{2k+1}, S_{2k}, T_{2k+1}$ in the same fashion, which results in the following system of equations:

\begin{equation} \label{eq:PRST}
\begin{cases}
\begin{aligned}
P_{2k} &= 1 + {\color{red}z C_{e,2k-2} S_{2k}} + {\color{blue}z C_{o,2k-1}P_{2k}}\\
R_{2k+1} &= \phantom{1+\ \ } {\color{blue}z C_{e,2k} T_{2k+1}} + {\color{red}z C_{o,2k-1} R_{2k+1}}\\
S_{2k} &= \phantom{1+\ \ } {\color{blue}z C_{e,2k-2} P_{2k}} + {\color{red}z R_{2k+1} S_{2k}}\\
T_{2k+1} &= 1 + {\color{red}z  P_{2k+2} R_{2k+1}} + {\color{blue}z C_{o,2k-1} T_{2k+1}}
\end{aligned}
\end{cases}
\end{equation}

%\vspace{0.25in}

Now, consider the sequence $\{|\mathfrak{D}^4_{2n}(1423)|\}_{n\ge 0}=\{|\mathfrak{D}^4_{2n}(312)|\}_{n\ge 0}$. We claim that its ordinary generating function is $R_1/z$. Indeed, removing the top row and the rightmost column of the $R$-board with $m=1$ yields exactly the board of allowed cells in a Dumont-4 permutation.

Now consider the function $R_1/z$. Solving the system of equations \eqref{eq:PRST}  for $R_{2k+1}$, $k\ge 0$, yields the following recursive formula:
\[
R_{2k+1} = \frac{zC_{e,2k}}{(1-zC_{o,2k-1})^2} \,\,
 \genfrac{}{}{0pt}{}{}{-} \,\,
\frac{z^2C_{e,2k}}{1-zC_{o,2k+1}} \,\,
 \genfrac{}{}{0pt}{}{}{-} \,\,
\frac{z^2C_{e,2k}^2}{1} \,\,
 \genfrac{}{}{0pt}{}{}{-} \,\, 
zR_{2k+3}.
\]

Multiplying both sides by $z$, we obtain for all $k\ge 0$:
\[
zR_{2k+1} = \frac{z^2C_{e,2k}}{(1-zC_{o,2k-1})^2} \,\,
 \genfrac{}{}{0pt}{}{}{-} \,\,
\frac{z^2C_{e,2k}}{1-zC_{o,2k+1}} \,\,
 \genfrac{}{}{0pt}{}{}{-} \,\,
\frac{z^2C_{e,2k}^2}{1} \,\,
 \genfrac{}{}{0pt}{}{}{-} \,\, 
zR_{2k+3}.
\]

Note, the term on the left and the last term on the right are of the same form with $k$ increasing by $1$. This yields the continued fraction representation for the generating function $R_1/z=zR_1/z^2=\sum_{n=0}^{\infty}{|\mathfrak{D}^4_{2n}(312)|z^{2n}}$. Substituting $\sqrt{z}$ for $z$, we get the generating function for $\{|\mathfrak{D}^4_{2n}(312)|\}_{n\ge 0}$.
\end{proof}

Also note that the numerators  of first three terms in the recurrence formula contain only the truncations of the even part of $C(z)$, and the denominators contain only the truncations of the odd part of $C(z)$. Moreover, all of those numerators and denominators are even functions. Furthermore, it is not difficult to see that the term $zR_{2k+3}$ in the recursive formula \eqref{eq:d4-312-ogf} does not contribute to the terms of degree at most 6 in $zR_{2k+1}$ (due to the factors of $z^2$ in the three numerators). Applying this observation iteratively, we see that deleting the term $zR_{2k+1}$ in the resulting recursive formula for $R_1/z$ does not affect the terms of degrees $2n\in[0,6k-2]$, i.e. $2n$ for $0\le n\le 3k-1$.

Finally, we note that the sequence $\{|\mathfrak{D}^4_{2n}(312)|\}_{n\ge 0}$ is \href{https://oeis.org/A343795}{A343795} in OEIS \cite{OEIS} and begins 
\[
1,1,3,10,39,174,872,4805,28474,178099,1160173,7803860,\dots .
\]
Note that these are the coefficients for $0\le n\le 11$, so $3k-1=11$ yields $k=4$, and therefore these terms can be found by expanding $R_1/z$ iteratively until we reach $zR_9$, then removing the term $zR_9$.

% End Section 3 - Avoidance on Dumont permutations of the fourth kind

%%%%%%%%%%%%%%%%%%%%

% Begin Section 4 - A single occurrence of patterns in Dumont permutations

\section{A single occurrence of patterns in Dumont permutations}

The results of this section focus on enumeration of Dumont permutations with a single occurrence of certain patterns.

\subsection{One occurrence in Dumont-1 and Dumont-2 permutations}

We first introduce notation that is useful for going back and forth between sequences and their generating functions. Then we will review results on single occurrences of patterns in Dumont-1 and Dumont-2 permutations. Finally, we will prove a related result on a single occurrence of a pattern in Dumont-4 permutations.

\begin{notation}
For any ordinary generating function, let:
\[
A(z) \longleftrightarrow \{a_n\} \quad \text{if} \quad A(z) = \text{ogf}\{a_n\} = \sum_{n=0}^\infty a_n z^n.
\]
\end{notation}

\noindent For example, for the Catalan numbers we have $\{C_n\} \longleftrightarrow  C = C(z)$ and $C = 1 +  zC^2$.
Recall that
\[
C_n = \frac{1}{n+1} \binom{2n}{n},
\]
which implies that, for any $k \ge 1$,
\[
C^k \longleftrightarrow \frac{k}{n+k} \binom{2n+k-1}{n}.
\]

%\vspace{0.2in}

%\noindent 
Likewise, for the central binomial coefficients $\displaystyle B_n = \binom{2n}{n}$ we have $\{B_n\} \longleftrightarrow  B = B(z)$ and $B = 1 +  2zBC.$  This implies that, for any $k \ge 1$, 
$$BC^k \longleftrightarrow \binom{2n+k}{n}.$$ 

\begin{theorem}[\cite{Mansour}] \label{thm:132;1} 
For all $n\ge 0$, there does not exist a Dumont-1 permutation containing $132$ exactly once.  That is,
\[
|\mathfrak{D}^1_{2n} (132;1)| = 0.
\]
\end{theorem}
%
%We reproduce the proof of this theorem to show how the block decompostion is used to determine the structure and enumerate permutations containing a single occurrence of a pattern, in this case $132$.
%
%\begin{proof}
%We will proceed by induction on $n$.  
%
%Let $\pi$ be a Dumont-1 permutation of length $2n$ that contains exactly one occurrence of $132$, where $\pi = (\pi', 2n-1, 2n, \pi'')$ or $\pi=(2n,\pi',2n-1)$, and $\pi''$ is nonempty. It is straightforward that there does not exist a Dumont-1 permutation for $n=0,1$ containing an occurrence of pattern $132$.  So, suppose $n \ge 2$, and assume by induction on $n$ that there does not exist a Dumont-1 permutation of length $k \le 2n-1$ containing exactly one occurrence of $132$.  To prove this property for $2n$, consider the following two cases: where (1) $2n$ does not appear in the single occurrence of $132$, or (2) $2n$ does appear in the single occurrence of $132$.  Before continuing, we need a proposition about Dumont-1 permutations avoiding 132 and the position of that maximal element $2n$.
%
%%\vspace{0.25in}
%\begin{prop}
%\label{132_prop}
%\cite{Mansour} 
%For any $\pi \in \mathfrak{D}^1_{2n} (132)$ such that $\pi_j = 2n$, we have 
%\begin{itemize}
%\item $\pi = (\pi'+(2n-2j), 2n-1, 2n, \pi'')$, where
%$\pi'\in\mathfrak{D}^1_{2j-2} (132)$,\\ $\pi''\in\mathfrak{D}^1_{2n-2j} (132)$ $(j = 1, 2, 3, \dots, n-1)$, and $\pi''$ is nonempty, or
%\item $\pi=(2n,\pi',2n-1)$, where $\pi'\in\mathfrak{D}^1_{2n-2} (132)$.
%\end{itemize}
%\end{prop}
%
%We now continue with the proof of Theorem \ref{thm:132;1}.
%
%\noindent
%\underline {Case 1:}
% Let $2n$ not appear in the single occurrence of  $\pi$.  Then $\pi$ has the same structure as in Proposition \ref{132_prop}, except that exactly one of $\pi'$ or $\pi''$ contains a single instance of $132$, while the other one avoids $132$ (see Figure \ref{fig:blocks132_single} on page \pageref{fig:blocks132_single}). But then, by induction on $n$, this is impossible, because Dumont-1 permutations of size smaller than $2n$ containing a single occurrence of 132 do not exist.
%
%\begin{figure}[!h]
%\begin{center}
%\begin{tikzpicture} [scale = 0.75]
%
%\begin{scope}
% \foreach \position in {(3.5,5.5), (4.5,6.5)}
%    \filldraw [color = black] \position circle (4pt);
% \node [right] at (3.75,5.5) {$2n-1$};
% \node [right] at (4.75,6.5) {$2n$};
%\draw [line width = 2pt, color=black] (0,2) rectangle (3,5);
% \node at (1.5,3.5){$\pi'$};
%\draw [line width = 2pt, color=black] (5,0) rectangle (7,2);
% \node at (6,1){$\pi''$};
%\draw [line width = 1pt, color=black, dashed] (-0.5,2) -- (7.5,2);
% \node at (8.25,2.1){$2n-2j$};
% \node at (3, 8) {\underline{$\pi''$ nonempty}};
%\end{scope}
%
%\begin{scope}[xshift=13cm]
% \foreach \position in {(6.5,5.5), (0.5,6.5)}
%    \filldraw [color = black] \position circle (4pt);
% \node [right] at (6.75,5.5) {$2n-1$};
% \node [right] at (0.75,6.5) {$2n$};
%\draw [line width = 2pt, color=black] (1,0) rectangle (6,5);
% \node at (3.5,2.5){$\pi'$};
%  \node at (3.5, 8) {\underline{$\pi''$ empty}};
%\end{scope}
%
%\end{tikzpicture} 
%\caption{Block decompositions for $\pi \in \mathfrak{D}^1_{2n} (132;1)$; $2n$ is not in the 132-occurrence}
%\label{fig:blocks132_single}
%\end{center}
%\end{figure}
%
%\noindent
%\underline {Case 2:}
%Let $2n$ appear in the single occurrence of $132$ of $\pi$.  Then by a similar analysis as in Proposition \ref{132_prop} (see Figure \ref{fig:block132_single} on page \pageref{fig:block132_single}), $\pi = (\pi', k, 2n, \pi'', k+1, \pi''')$ such that $\pi_p = 2n$ and $\pi_q = k+1,$ where every element of $\pi'$ is greater than every element of $\pi''$, and every element of $\pi''$ is greater than every element of $\pi'''$. 
%
%\begin{figure}[!h]
%\begin{center}
%\begin{tikzpicture} [scale = 1]
%
%% 1st column
%\draw [line width = 1pt, color=black, fill=lightgray] (0,0) rectangle (1,2);
%\draw [line width = 1pt, color=black, fill=lightgray] (0,2) rectangle (1,3);
%\draw [line width = 1pt, color=black] (0,3) rectangle (1,4);
%% 2nd column
%\draw [line width = 1pt, color=black, fill=lightgray] (1,0) rectangle (2,2);
%\draw [line width = 1pt, color=black, fill=lightgray] (1,2) rectangle (2,3);
%\draw [line width = 1pt, color=black, fill=lightgray] (1,3) rectangle (2,4);
%% 3rd column
%\draw [line width = 1pt, color=black, fill=lightgray] (2,0) rectangle (3,1);
%\draw [line width = 1pt, color=black] (2,1) rectangle (3,2);
%\draw [line width = 1pt, color=black, fill=lightgray] (2,2) rectangle (3,3);
%\draw [line width = 1pt, color=black, fill=lightgray] (2,3) rectangle (3,4);
%% 4th column
%\draw [line width = 1pt, color=black] (3,0) rectangle (4,1);
%\draw [line width = 1pt, color=black, fill=lightgray] (3,1) rectangle (4,2);
%\draw [line width = 1pt, color=black, fill=lightgray] (3,2) rectangle (4,3);
%\draw [line width = 1pt, color=black, fill=lightgray] (3,3) rectangle (4,4);
%% dots
%\foreach \position in {(1,2), (2,4), (3,3)}
%    \filldraw [color = black] \position circle (3pt);
%% nodes
%\node at (0.5,3.5) {$\pi'$};
%\node at (2.5,1.5) {$\pi''$};
%\node at (3.5,0.5) {$\pi'''$};
%\node at (1.25,2.25) {$k$};
%\node at (2.35,4.25) {$2n$};
%\node at (3.5,3.25) {$k\!+\!1$};
%
%\end{tikzpicture} 
%\caption{Block decompositions for $\pi \in \mathfrak{D}^1_{2n} (132;1)$; $2n$ appears in the 132-occurrence}
%\label{fig:block132_single}
%\end{center}
%\end{figure}
%
%It follows, since $2n$ is even and maximal, then the preceding entry $k$ must be odd, thus $k+1$ is even. Therefore, $\pi'''$ must be nonempty. Note that in a Dumont-1 permutation, the entry $2n-1$ must appear either immediately before $2n$ or at the end. However, the entry $k$ immediately before $2n$ is at most $2n-2$, and the last entry of $\pi$ is the last entry of $\pi'''$, which is less than $k$, so also cannot be $2n-1$.
%%
%Therefore, by induction on $n$, our theorem is proved.
%\end{proof}

%The next few theorems have proofs which use similar methods, thus will be introduced without proofs. 

The next few theorems list previous results \cite{Burstein_Ofodile, Ofodile} on Dumont-1 and Dumont-2 permutations with a single occurrence of certain patterns.

%\vspace{0.1in}
\begin{theorem}[\cite {Ofodile}]
For all $n\ge 0$, 
\[
|\mathfrak{D}^1_{2n} (312;1)| = 0 \quad \text{ and } \quad |\mathfrak{D}^1_{2n} (231;1)| = \binom {2n-2} {n-3}.
\]
\end{theorem}

%\begin{theorem}
%\cite {Ofodile}
%For all $n\ge 0$, there does not exist a Dumont-1 permutation containing $312$ exactly once.  That is,
%$$|\mathfrak{D}^1_{2n} (312;1)| = 0.$$
%\end{theorem}
%
%%\vspace{0.1in}
%\begin{theorem}
%\cite {Ofodile}
%For all $n\ge 3$,
%$$|\mathfrak{D}^1_{2n} (231;1)| = \binom {2n-2} {n-3}.$$
%\end{theorem}

%\noindent
Note that 
%since $n-3$ is the number of elements we wish to ``choose" from $2n-2$, isolating an ``$n-3$" from the set that we are choosing from, we have 
$2n-2=2(n-3)+4$, so
\[
\binom {2n-2} {n-3} = \binom {2(n-3)+4} {n-3} \longleftrightarrow z^3BC^4.
\]

%\vspace{0.1in}
\begin{theorem}[\cite {Ofodile}]
For all $n\ge 4$,
\[
|\mathfrak{D}^1_{2n} (213;1)| = C_{n-2}+\binom {2n-4} {n-4}.
\]
\end{theorem}

%\noindent
Similarly,
\[
C_{n-2}+\binom {2n-4} {n-4} = C_{n-2}+\binom {2(n-4)+4} {n-4} \longleftrightarrow z^{2}C+z^4BC^4.
\]

%\vspace{0.1in}
\begin{theorem}[\cite{Burstein_Ofodile}]
For all $n \ge 2$,
\[
|\mathfrak{D}^1_{2n} (321;1)| = (n-1)^2.
\]
\end{theorem}

%\vspace{0.1in}
\begin{theorem}[\cite{Burstein_Ofodile}]
For all $n \ge 2$,
\[
|\mathfrak{D}^2_{2n} (321;1)| = \frac{5}{n+3}\binom{2n}{n-2} \longleftrightarrow z^2 C^5.
\]
\end{theorem}

%\vspace{0.1in}
\begin{theorem}[\cite {Ofodile}]
For all $n\ge 2$,
\[
|\mathfrak{D}^2_{2n} (3142;1)| = \binom {2n-1} {n-2} \longleftrightarrow z^2BC^3.
\]
\end{theorem}

%\vspace{0.1in}
\begin{theorem}[\cite{Burstein_Ofodile}]
For all $n \ge 2$,
\[
|\mathfrak{D}^2_{2n} (2143;1)| = a_n b_{n+1} + b_n a_{n+1} + a_{n-1}a_n,
\]
where
\[
a_{2k} = \frac{1}{2k+1} \binom{3k}{k}, \quad a_{2k+1} = \frac{1}{k+1} \binom{3k+1}{k},
\]
and
\[
b_{2k} = \binom{3k-3}{k-2}, \quad b_{2k+1} = 2 \binom {3k-2}{k-2}.
\]
\end{theorem}

\subsection {A single occurrence of patterns in Dumont-4 permutations}

The results of this subsection focus on a single occurrence of patterns in Dumont-4 permutations.  To begin, let us consider a theorem.

\begin{theorem}[\cite {Noonan, Noonan_Zeilberger}] \label{Noo_Zeil}
\[
|S_n(123;1)| = \frac{3}{n} \binom{2n}{n-3}
\]
%\label{Noonan_Thm}
\end{theorem}
%
%Theorem \ref {Noo_Zeil} which states that the number of permutations with one occurrence of the pattern $123$ is $\displaystyle \frac{3}{n} \binom{2n}{n-3}.$  
% 
This result enumerates all permutations containing a single copy of pattern $321$. The original proof by both Noonan \cite {Noonan} and Noonan and Zeilberger \cite {Noonan_Zeilberger} uses a complicated induction and generating functions with multiple auxiliary parameters. A much shorter argument has been subsequently given by the first author in \cite{Burstein_Short}, which was then shortered further by Zeilberger \cite{Zeilberger_lovely} and recently generalized by B\'ona and the first author in \cite{Bona_Burstein}.  
%Additionally, Burstein proved it using the reversal $321$ of $123$ for an additional nice property.  
Given the methods used in the proofs in \cite{Burstein_Short, Zeilberger_lovely}, as well as the corresponding permutation diagram of the decomposition of permutations in $S_n(321;1)$ as in \cite{Burstein_Short}, we will now prove a theorem first presented in \cite{Burstein_Jones}, enumerating a single occurrence of $321$ in Dumont-4 permutations.

%
%Additionally, Burstein proved it using the reversal $321$ of $123$ for an additional nice property.  We reproduce the ``lovely" proof by Burstein and not the ``shortened lovely" proof by Zeilberger, that $\displaystyle |S_n(123;1)| = \frac{3}{n} \binom{2n}{n-3}$, to show a necessary injection for the results of this section.
%%
%\begin{proof}
%\cite {Burstein_Short}
%Let a permutation, say $\pi$, of length $n$ have exactly one occurrence of the pattern $321$.  We can write $\pi$ as $\pi_1 c \pi_2 b \pi_3 a \pi_4$, where the subsequence $(c,b,a)$ is the one and only unique occurrence of $321$, with $a < b < c$.  Also, let the occurrence of $321$ in $\pi$ be at positions $(j, k, l)$, with $j < k < l$. Note, all entries to the left of $b$ must be smaller than $b$, with the exception of $c$.  Likewise, all entries to the right of $b$ must be larger than $b$, with the exception of $a$.  If this were not the case, we would indeed have another occurrence of $321$.  Thus, if we were to exchange entries $c$ and $a$ at positions $j$ and $l$, that is taking our one occurrence of $321$ and turning it into an occurrence of $123$ (hence avoiding $321$), we must have the position $k$ and the entry $b$ as equal.  In other words, $k=b$ results in $b$ as a fixed point of $\pi$.  Note, this is the added nice property Burstein \cite{Burstein_Short} referenced by considering the reverse of $123$ instead of $123$ itself.
%
%Now, we define a map $f: S_n(321;1) \longrightarrow S_{n+2}(321)$.  Again $\pi = \pi_1 c \pi_2 b \pi_3 a \pi_4$.  Consider a permutation diagram of the single occurrence of $321$.  Using block decomposition, we will replace each of the points $(j,c)$ and $(l,a)$ with two additional points as follows. For some $\epsilon$ such that $0 < \epsilon < 1$:
%
%\begin{itemize}
%\item Replace the point $(j,c)$ with two points $(j, b - \epsilon)$ and $(b + \epsilon, c)$.
%\item Replace the point $(l, a)$ with two points $(b - \epsilon, a)$ and $(l, b + \epsilon)$.
%\end{itemize}
%
%So, instead the subsequence $(c,b,a)$ which is the one occurrence of $321$ in $\pi$, we now have a permutation, say $\hat{\pi}$, with the subsequence $(a,b,c)$ and two new entries; one between $\pi_1$ and $\pi_2$, and the other  between $\pi_3$ and $\pi_4.$  That is,
%$$\hat{\pi} = \pi_1 (b-\epsilon) \pi_2 a b c \pi_3 (b+\epsilon) \pi_4.$$ 
%
%Thus, we have $\hat{\pi} = \pi' b \pi'',$ where:
%\begin{itemize}
%\item $\pi' = \pi_1 (b-\epsilon) \pi_2 a$ with entries from the set $\{1, 2, \dots, b-1, b-\epsilon \}$ and 
%\item $\pi'' = c \pi_3 (b+\epsilon) \pi_4$ with entries from the set $\{b+\epsilon, b+1, \dots, n-1, n\}$. 
%\end{itemize}
%
%We see that every entry in $\pi'$ is less than $b$, and every entry in $\pi''$ is greater than $b$, that is $\pi' < b < \pi''$.  Finally, define $f(\pi)$ to be the unique permutation in $S_{n+2}$ order-isomorphic to $\hat{\pi}$.
%
%\begin{figure}[!h]
%\begin{center}
%\begin{tikzpicture} [scale = 1.25]
%\fill[gray!50!white] (0.1,2) rectangle (2,3.9);
%\fill[gray!50!white] (2,0.1) rectangle (3.9,2);
%\fill[gray!50!white] (2,2) rectangle (3,3);
%\fill[gray!50!white] (1,1) rectangle (2,2);
%\draw[thin,color = black] (0.1,0.1) grid (3.9,3.9);
%
% \foreach \position in {(1,3), (2,2), (3,1)}
%    \filldraw [color = black] \position circle (4pt);
% \foreach \position in {(1,1.85), (2.15, 3), (1.85, 1), (3,2.15)} 
%    \filldraw [color = blue] \position circle (4pt);
%\draw [line width = 2pt, color = blue, dashed] (1.85, 0.1) -- (1.85,1.85) -- (0.1,1.85);
%\draw [line width = 2pt, color = blue, dashed] (2.15, 3.9) -- (2.15, 2.15) -- (3.9, 2.15);
%\draw [->, line width = 2.5pt, color=red] (1,2.85) -- (1,2);
%\draw [->, line width = 2.5pt, color=red] (1.15,3) -- (2, 3);
%\draw [->, line width = 2.5pt, color=red] (3, 1.15) -- (3, 2);
%\draw [->, line width = 2.5pt, color=red] (2.85,1) -- (2, 1);
%\end{tikzpicture} 
%	\caption {Permutation diagram of an injection 
%				$f:S_n(321;1) \longrightarrow S_{n+2}(321).$}
%	\label{fig:321-occurrence}
%\end{center}
%\end{figure}
%
%Now, it is straightforward that both $\pi'$ and $\pi''$ avoid $321$, therefore $f(\pi)$ also avoids $321$.  In other words $f(\pi) \in S_{n+2} (321)$.  Additionally, $f$ is an injection.  With respect to restrictions on $\pi'$ and $\pi''$, both are nonempty and the greatest (resp. least) entry of $\pi'$ (resp. $\pi''$) is distinct from its rightmost (resp. leftmost) entry.  Thus, there is a bijection between $S_n(321)$ and its image, that is the subset of $S_{n+2}(321)$ which consists of permutations $(\rho', b+1, \rho'')$ for $1 \le b \le n$ such that:
%\begin{itemize}
%\item $\rho'$ avoids $321$, has entries from the set $[b]$, and does not end on $b$, and
%\item $\rho''$ avoids $321$, has entries from the set $\{b+2, b+3, \dots, n+2\}$, and does not start with $b+2$.
%\end{itemize}
%
%So from the fact that $C = 1 + zC^2$ for the Catalan generating function $C$, we obtain that the generating function for the number of nonempty 321-avoiding permutations whose greatest (resp. least) value is not their rightmost (resp. leftmost) value is $$C - 1 - zC = zC^2 - zC = zC(C-1) = zC(zC^2) = z^2 C^3.$$  Hence, the generating function for the number of permutations with exactly one occurrence of $321$ is $$(z^2 C^3)^2 \cdot z^{-2} \cdot z = z^3 C^6$$ and therefore, the number of such permutations is 
%$\displaystyle \frac{3}{n} \binom {2n}{n-3}$.
%\end{proof}
%
%\noindent
%Note the following:
%\begin{itemize}
%\item the two factors $z^2 C^3$ correspond to $\pi'$ and $\pi''$, 
%\item the factor of $z^{-2}$ corresponds to the number $-2$ from $|\pi| - |f(\pi)| = n - (n+2) = -2,$ and
%\item the factor of $z$ corresponds to the fixed point $(b,b)$.
%\end{itemize}
%
%Given the methods used in the ``lovely" proof of \cite{Burstein_Short}, as well as the permutation diagram of the injection $f:S_n(321;1) \longrightarrow S_{n+2}(321),$ we will now prove a theorem first presented in a talk \cite{Burstein_Jones}, enumerating a single occurrence of $321$ in Dumont-4 permutations.

\begin{theorem}[\cite{Burstein_Jones}] \label{thm:d4-321-1}
For all $n\ge 1$,
\[
\begin{aligned}
|\mathfrak{D}^4_{2n} (321;1)| &= |S_n (321;1)| + |S_{n+1} (321;1)|\\
&= \frac{3}{n}\binom{2n}{n-3}+\frac{3}{n+1}\binom{2n+2}{n-2}\\
&= C_{n+3} - 3C_{n+2} - C_{n+1} + 3C_n.
\end{aligned}
\]
\end{theorem}

It follows from Theorems \ref{Noo_Zeil} and \ref{thm:d4-321-1} that the generating function for the sequence $|\mathfrak{D}^4_{2n}(321;1)|$, $n\ge 0$, is
\[
z^3 C^6 + \frac{1}{z}\cdot z^3 C^6 = (z^2+z^3)C^6.
\]

%\vspace {0.25in}

%\begin{proof}
%Considering that $\mathfrak{D}^4_{2n} (321;1) \subseteq S_{2n}(321;1)$, the ``2" is fixed and the ``3" and ``1" have similar properties if we look at moving each to their two respective locations to obtain 321-avoiding permutations which may or may not necessarily be Dumont-4 permutations.  We start by temporarily removing the first entry since by definition it is a 1, and we will also remove the ``2" of the occurrence of $321$ temporarily.  We then analyze two cases of the form of that particular permutation with the ``2" removed. Next, we add entries to our two permutations so as to result in two new Dumont-4 permutations that avoid 321, which as we know are enumerated by the Catalan numbers.  We view these two new Dumont-4 permutations as an ordered pair.  We will see that the combined length of the two resulting permutations in $\mathfrak{D}^4(321)$ is either $2n+2$ or $2n+4$, depending on whether the ``2" is even or odd, respectively.  We will use the map of Theorem \ref{thm:1432} to produce two  permutations that avoid 321 as in the proof in \cite {Burstein_Short} with the combined length of $n+1$ or $n+2$.  We then combine them to produce a single Dumont-4 permutation of length $n$ or $n+1$ (depending on the parity of ``2'') with a single occurrence of pattern 321.

%\begin{proof}

%To begin, we have a Dumont-4 permutation with only one occurrence of the pattern 321.  We take note of the entry that represents the ``2" in the 321-occurrence, whether it is even or odd.  And again, the ``2" is a fixed point, so if the entry is even it will occur in an even position, likewise for an odd entry. 

\begin{proof}
Let $\pi \in \mathfrak{D}^4_{2n}(321;1)$. Suppose that the single occurrence of pattern $321$ is formed by values $c>b>a$ and $\pi=\tau_1c\tau_2b\tau_3a\tau_4$ for some strings $\tau_i$, $i=1,2,3,4$. Then $b$ is a fixed point. Let $\pi'$ and $\pi''$ be the patterns of $\tau_1c\tau_2a$ and $c\tau_3a\tau_4$, respectively, as in \cite{Zeilberger_lovely}.  Then $\pi'$ and $\pi''$ avoid $321$, $\pi'$ does not end on its top entry, and $\pi''$ does not start with its bottom entry. Moreover, the entry $a$, the ``1'' of the occurrence of $321$, is a deficiency and hence an even value in an even position in $\pi$. 

Consider two cases based on the parity of $b$, the entry ``$2$'' of the single $321$-occurrence in $\pi$.

\begin{enumerate}[\emph{Case} 1:]

\item \label{enum:b-even} Suppose $b$ is even, say, $b=2k$ for some $k\in[2,n-1]$ (note that the entry $b$ cannot be either $2$ or $2n$, since neither is involved as a ``2'' in an occurrence of $321$: the former because $1$ is a fixed point, the latter because $2n$ is the largest entry).

Then $|\pi'|=2k$ and $|\pi''|=2n-2k+1$. Recall that $a$ is even; moreover, it occupies the last position in $\pi'$, i.e. position $2k$, which is even. Therefore, $\pi'$ is also Dumont-4, and thus $\pi'\in\mathfrak{D}^4_{2k}(321)$. For what follows after this case analysis, let $\rho'=\pi'$.

Now consider the permutation $\rho''=(1,\pi''+1)$.  Note that the entry 2 in $\rho''$ is not a fixed point, since $\pi''$ does not begin with its lowest entry, thus 2 is a deficiency of $\rho''$. Therefore, $\rho''$ is also Dumont-4, and thus $\rho''\in\mathfrak{D}^4_{2n-2k+2}(321)$.

Note that the sizes of $\rho'$ and $\rho''$ add up to $2k+(2n-2k+2)=2n+2$ in this case.

\item \label{enum:b-odd} Suppose $b$ is odd, say, $b=2k+1$ for some $k\in[1,n-1]$.  Then $|\pi'|=2k+1$ and $|\pi''|=2n-2k$.  

Insert the new entry $2k+2$ in the second rightmost position of $\pi'$ to form a permutation $\rho'$.  Then $|\rho'|=2k+2$, which is even. Moreover, as in Case \ref{enum:b-even}, the entry $a$, which is at most $2k$ and even (as a deficiency), now occupies position $2k+2$, also even. Therefore, $\rho'\in\mathfrak{D}^4_{2k+2}(321)$.

Now form the permutation $\rho''$ from $\pi''$ in the following fashion: if $\pi''=(\pi_1, 1, \pi_2)$, then $\rho''=(1,3,\pi_1+3,2,\pi_2+3)$. Note that $|\rho''|=|\pi''|+2=2n-2k+2$ all deficiencies of $\rho''$ come from deficiencies of $\pi$. 

The entry $a$ in $\pi$ becomes the entry $2$ of $\rho''$, and hence also an even deficiency. Its position in $\rho''$ is $(\rho'')^{-1}(2)=\pi^{-1}(a)-(2k+1)+1+2=\pi^{-1}(a)-(2k-2)$ is also even, since $\pi^{-1}(a)$ is even. The entries of $\rho''$ that are at least $4$ are obtained by subtracting $2k-2$ from the corresponding entries of $\pi$. However, their position also shifted to the left by $2k-2$, so the parity of their positions and values as well as the placement relative to the diagonal (i.e. being a fixed point, excedance, or deficiency) remains the same. Therefore, $\rho''$ is a Dumont-4 permutation, and thus $\rho''\in\mathfrak{D}_{2n-2k+2}(321)$.

Note that the sizes of $\rho'$ and $\rho''$ add up to $(2k+2)+(2n-2k+2)=2n+4$ in this case.

\end{enumerate}

%\bigskip

Now, given permutations $\rho'$ and $\rho''$, use the map of Theorem \ref{thm:1432} to produce two permutations $\sigma'$ and $\sigma''$, respectively, that avoid $321$. Note that $\rho'$ does not end on its top entry, and therefore neither does $\sigma'$. Likewise, $\rho''$ does not start with $12\dots$, and therefore $\sigma''$ does not start with $1$. Thus, the pair $\sigma'$ and $\sigma''$ are as in the proof in \cite {Burstein_Short} with the combined length of $(2n+2)/2=n+1$ or $(2n+4)/2=n+2$. Thus, we can combine them as in \cite {Burstein_Short} to produce a single permutation $\sigma$ of length $n$ or $n+1$ (depending on the parity of $b$) with a single occurrence of pattern $321$.

%\hrulefill
%
%Next, we remove the ``2" in the occurrence of 321, which is a fixed point and again reduce the entries of the resulting permutation accordingly. This will allow us to analyze two split substrings by capturing all entries below (resp. above) the removed ``2", along with the ``3" (resp. ``1") which will be the only entry above (resp. below) the removed ``2".  Now with the removed ``2", we have two substrings that both avoid 321.  Their lengths are as follows:
%
%\begin{itemize}
%\item if the ``2" was originally even, the two substrings are of length $2k+1$ and $2n-(2k+1)$ for some $1\le k\le n-1$, and
%\item if the ``2" was originally odd, the two substrings are of length $2k$ and $2n-2k$ for some $1\le k\le n-1$.
%\end{itemize}
%
%\noindent
%Now to reinsert the temporarily removed entries, we do so under the following conditions:
%\begin{itemize}
%\item If the ``2" was originally even, for the two substrings that are of length $2k+1$ and $2n-(2k+1)$, replace the 1 in the beginning of both of the substrings.
%\item If the ``2" was originally odd, for the two substrings that are of length $2k$ and $2n-2k$, 
%\begin{itemize}
%\item for the substring that was formed from capturing all entries below the removed ``2" along with the ``3", insert the maximal entry in the second to last position, and then replace the 1 to the beginning, and
%\item for the substring that was formed from capturing all entries above the removed ``2" along with the ``1", insert the second minimal entry at the beginning, and then replace the 1 to the very beginning.
%\end{itemize} 
%\end{itemize} 
%
%\noindent
%This algorithm preserves the number of right-to-left minima from the substrings, and presents an ordered pair of Dumont-4 permutations that avoid 321.  Before continuing, we show an example of a Dumont-4 permutation with a single occurrence of 321, which will result in an ordered pair of Dumont-4 permutations that avoid 321.
%
%
%Continuing with our proof, we see that our Dyck paths from both of the Dumont-4 permutations in the ordered pair, after collapsing them in half from semilength $2j$ to semilength $j$ (see Subsection \ref{321Dyck}, specifically Figure \ref{fig:Dyck6x6_1432}) can be placed on a permutation diagram grid in their original orientation. Here, the Dyck paths will not need the extra space below the right-to-left minima as mentioned in the proof of Theorem \ref{321thm}.  After placing the right-to-left minima where a corner arises from the Dyck path, any other spaces will be filled in an increasing manner.  Note, these permutations still avoid 321; if the remaining spaces were not filled in with dots in an increasing manner we would have one or more occurrences of 321 .  
%
%Lastly, to complete the injective map we recreate the single occurrence of 321 by analyzing the common elements from the ``3" and the ``1" from both of the newly formed permutations from collapsing the Dyck paths of semilength $2j$ to semilength $j$.  We now replace the removed ``2" between the ``3" and the ``1" in the one and only position that allows for a single occurrence of 321:
%\begin{itemize}
%\item for the substring that was captured by choosing the lower entries, along with the ``3", replace the ``2" between the ``3" and the ``1" in the rightmost location, and
%\item for the substring that was captured by choosing the upper entries, along with the ``1", replace the ``2" between the ``3" and the ``1" in the leftmost location.
%\end{itemize}
%
%One more note is that the result that will yield the desired permutation will be that of greatest length from the pair of substrings after the Dyck paths were collapsed from semilength $2j$ to semilength $j$; when the original ``2" was even, we had two substrings both of length $n-1$, so adding the removed ``2" gives us a permutation of length $n$ with one occurrence of 321.  Likewise, when the original ``2" was odd, we had two substrings of length $n-1$ and length $n$, so adding the removed ``2" gives us a permutation of length $n+1$ with one occurrence of 321.

Therefore, the number of Dumont-4 permutations of length $2n$ that have a single 321-occurrence is equal to the number of permutations of length $n$ with a single 321-occurrence plus the number of permutations of length $n+1$ with a single 321-occurrence.  That is, $$|\mathfrak{D}^4_{2n} (321;1)| = |S_n (321;1)| + |S_{n+1} (321;1)|.$$  Finally, Zeilberger \cite {Zeilberger_lovely} showed that $|S_n (321;1)| = C_{n+2} - 4C_{n+1} + 3C_n$.  It follows that 
\[
%\begin{split}
|\mathfrak{D}^4_{2n} (321;1)| =( C_{n+2} - 4C_{n+1} + 3C_n)+ ( C_{n+3} - 4C_{n+2} + 3C_{n+1})%\\
= C_{n+3} - 3C_{n+2} - C_{n+1} + 3C_n. \qedhere
%\end{split}
\]
\end{proof}

\begin{example} 
Let $\pi = 135462 \in \mathfrak{D}^4_{6}(321;1)$.  Here, we see that the occurrence of $321$ is the subsequence $542$, where the ``$2$" is the entry $4$, an even fixed point. The map of the proof in \cite{Burstein_Short} yields $\pi'=1342$ and $\pi''=231$, so that $\rho'=1342$ and $\rho''=1342$ as well. Now applying the map of Theorem \ref{thm:1432} yields $\sigma'=21$ and $\sigma''=21$. Combining those again as in \cite{Burstein_Short} yields $\sigma=321$. Note that $|\sigma|=|\pi|/2$.

Similarly, let $\pi = 136254 \in \mathfrak{D}^4_{6}(321;1)$.  Here, we see that the occurrence of $321$ is the subsequence $654$, where the ``$2$" is the entry $5$, an odd fixed point. The map of the proof in \cite{Burstein_Short} yields $\pi'=13524$ and $\pi''=21$, so that $\rho'=135624$ and $\rho''=1342$. Now applying the map of Theorem \ref{thm:1432} yields $\sigma'=312$ and $\sigma''=21$. Combining those again as in \cite{Burstein_Short} yields $\sigma=4132$. Note that $|\sigma|=|\pi|/2+1$.
%
%Now, removing the 1 in the beginning yields a permutation of length 5, which can be renamed as 25143.  Now, the entry 4 is the ``2", again a fixed point.
%
%At this point, we remove the ``2", capture all entries below (above, resp.) the removed ``2" along with the ``3" (``1", resp.) for the first (second, resp.) substring. Again, renaming the entries yield the ordered pair (2413, 21).  These substrings now avoid 321. To keep the avoidance of 321 and get back to a Dumont-4 permutation, 
%\begin{itemize}
%\item for the substring 2413 we add the maximal entry to the left of the entry 3, and a 1 to the beginning, which yields 135264 after renaming, and
%\item for the substring 21 we add the second minimal entry to the beginning, and then a 1 to the very beginning, which yields 1342 after renaming.
%\end{itemize}
%
%\noindent
%We now have the ordered pair $(135264, 1342)$ which consists of Dumont-4 permutations that avoid 321.
%\end{example} 
%
%
%\begin{example} 
%Continuing from our most previous example, we saw from $\pi = 136254 \in \mathfrak{D}^4_{6}(321;1)$, we obtained an ordered pair of Dumont-4 permutations which both avoid 321, that is (135264, 1342).
%
%Now, from analyzing their respective Dyck paths, and collapsing them from semilength 12 to semilength 6, reorienting those new Dyck paths on a permutation grid, placing the right-to-left minima along with filling in spaces, we arrive at the ordered pair (312, 21).  Lastly, placing the removed ``2" in its rightful place and renaming, we get the desired permutation of length 4 with a single occurence of 321. That is $\rho = 4132 \in S_4(321;1).$
\end{example} 


% End Section 4 - A single occurrence of patterns in Dumont permutations

%%%%%%%%%%%%%%%%%%%%

% Begin Section 5 - Conjecture and conclusion

\section{Conjecture and conclusion}

\subsection {A conjecture on restricted Dumont-1 permutations}
\label {B-J Conjecture}

The table at the end of \cite[Section 3.1]{Burstein_Restricted} implies that no two Dumont-1 permutations of length $4$ (recall: $\mathfrak{D}^1_4 = \{2143, 3421, 4213\}$) are Wilf-equivalent.  However, we found that this is due to an incorrect calculation of $|\mathfrak{D}_{10}(3421)|$, which is, in fact, equal to $|\mathfrak{D}_{10}(2143)|$. Furthermore, we found that $|\mathfrak{D}_{10}(3421)|=|\mathfrak{D}_{10}(2143)|$ for $n \le 6$; an additional computation by Albert \cite{Albert_comm} showed that the two sequences are the same for $n\le 10$, i.e. for Dumont-1 permutations of length up to $20$.  Therefore, we conjecture that the two sequences are equinumerous for all $n$.

\begin{conjecture}[\cite{Burstein_Jones}]
The following Wilf-equivalence holds on Dumont-1 permutations:
\[
2143 \sim 3421 \quad \text{on} \quad \mathfrak{D}^1_{2n}, \quad \text{that is} \quad |\mathfrak{D}^1_{2n} (2143)| = |\mathfrak{D}^1_{2n} (3421)| \quad \text{for all $n\ge 0$}.
\]
\end{conjecture}

Our conjecture is based on the following empirical observation. In all cases where sequences $|S_n(\pi_1)|$ and $|S_n(\pi_2)|$ are known to diverge (where $\pi_1$ and $\pi_2$ are patterns of the same length $k$), the first index at which the corresponding terms differ was found to be at most $2k-1$. Similarly, we expect Dumont permutations avoiding patterns of length $k$ to be Wilf-equivalent if their avoidance sequences coincide for avoiding permutations of length up to $4k-2$. In our case, the patterns are of length $k=4$, and the sequences coincide up to $n=10$, i.e. up to length $20$, greater than $4k-2=14$. Thus, we conjecture that they coincide for all $n\ge 0$.  See Table \ref{table:D1} for the sequence of the enumeration of $|\mathfrak{D}^1_{2n} (2143)| = |\mathfrak{D}^1_{2n} (3421)|$ up to $n = 10$.

%\vspace{0.25in}
\begin{table}[ht]
\begin{center}
\begin{tabular}{c||c|c|c|c|c|c|c|c|c|c|c}
	$n$ &0 & 1 & 2 & 3 & 4 & 5 & 6 & 7 & 8 &9 &10\\
	\hline 
	$|\mathfrak{D}^1_{2n} (2143)|$ & 1 & 1 & 2 & 7 & 36 & 239 & 1,892 & 17,015 & 168,503 & 1,799,272 & 20,409,644\\
	\hline 
	$|\mathfrak{D}^1_{2n} (3421)|$ & 1 & 1 & 2 & 7 & 36 & 239 & 1,892 & 17,015 & 168,503 & 1,799,272 & 20,409,644
\end{tabular}
\caption{Sequence of the cardinalities $|\mathfrak{D}^1_{2n} (2143)| = |\mathfrak{D}^1_{2n} (3421)|$}
\label{table:D1}
\end{center}
\end{table}

%\vspace{0.25in}

This conjecture has turned out to be very difficult to prove, and only some partial inroads are made as of this writing. One approach to proving such a conjecture is to refine it by considering distributions of some combinatorial statistics on both $\mathfrak{D}^1_{2n} (2143)$ and $\mathfrak{D}^1_{2n} (3421)$. For this, we will first need to define a generalized version of a permutation pattern.

\begin{definition}[\cite{Babson}]
A \emph{vincular} (or \emph{generalized}, or \emph{dashed}) permutation pattern is a pattern that allows the additional requirement that two (or more) adjacent letters in a pattern be also adjacent in the containing permutation.
\end{definition}

Note, in the pattern $2\text{-}31$, the ``$2$'' and ``$3$'' need not be adjacent, however the ``$3$'' and the ``$1$'' must be adjacent.  Now that we have the definition for vincular patterns, we can refine our conjecture slightly.  After studying the behavior of distributions of occurrences of various vincular patterns on $\mathfrak{D}^1_{2n}(2143)$ and $\mathfrak{D}^1_{2n}(3421)$, we were able to form the following conjecture.

\begin{conjecture}[\cite{Burstein_Jones}]
For all $n \ge 0$, we conjecture the following.
Let
\[
%\begin{split}
a_{n,k} = |\{\pi \in \mathfrak{D}^1_{2n} (2143) | (2\text{-}31)\pi = k\}|, \quad
b_{n,k} = |\{\pi \in \mathfrak{D}^1_{2n} (3421) | (13\text{-}2)\pi = k\}|,
%\end{split}
\]
where $(2\text{-}31)\pi$ (resp. $(13\text{-}2)\pi$) is the number of occurrences of $2\text{-}31$ (resp. $13\text{-}2$) in $\pi$. Then $a_{n,k} = b_{n,k}=1$ for $k=\binom{n}{2}$, $a_{n,k} = b_{n,k}=0$ for $k>\binom{n}{2}$, and
\[
\sum_{k=0}^{m}a_{n,k} \ge \sum_{k=0}^{m} b_{n,k} \quad \text{ for } 0 \le m \le \binom{n}{2},
\]
with equality for $m = \binom{n}{2}$.
\end{conjecture}

This is less optimal than statistic distributions that are equal to each other; however, of the statistics we studied, this relation is the only interesting one we found. Additionally, we conjecture that both sequences $(a_k)$ and $(b_k)$ are unimodal for each $n\ge 0$, equal at the tails of the distributions, while in the middle of the distributions we first have a block of $a_k > b_k$ followed by a block of $a_k < b_k$ with the switch in direction of the inequality occurring at approximately $k=2n-5$.  See Table \vref{table:cum_dist1} for $n=5$, Table \vref{table:cum_dist2} for $n=6$, and Table \vref{table:cum_dist3} for $n=7$.

%\vspace{0.25in}

\begin{table}[!h]
\begin{center}
%{\small
\begin{tabular}{c||c|c|c|c|c|c|c|c|c|c|c||c}
	$k$  &0 & 1 & 2 & 3 & 4 & 5 & 6 & 7 & 8 &9 &10 & \textbf{Total}\\
	\hline 
	$a_{5,k}$ & 1 & 10 & 30 & 45 & 49 & 42 & 31 & 18 & 9 & 3 & 1 & \textbf{239}\\
	\hline
	$b_{5,k}$ & 1 & 10 & 29 & 44 & 48 & 43 & 32 & 19 & 9 & 3 & 1 & \textbf{239}\\
	\hline
	$a_{5,k} \, \openbox \, b_{5,k}$ & $=$ & $=$ & $>$ & $>$ & $>$ & $<$ & $<$ & $<$ & $=$ & $=$ & $=$
\end{tabular}
%}
\caption{$|\{\pi \in \mathfrak{D}^1_{10} (2143) | (2\text{-}31)\pi = k\}|$ versus $|\{\pi \in \mathfrak{D}^1_{10} (3421) | (13\text{-}2)\pi = k\}|$}
\label{table:cum_dist1}
\end{center}
\end{table}

%\newpage
%\vspace{0.1in}

\begin{table}[!h]
%{\small
\begin{flushleft}
\begin{tabular}{c||c|c|c|c|c|c|c}
	$k$  &0 & 1 & 2 & 3 & 4 & 5 & 6 \\
	\hline 
	$a_{6,k}$ & 1 & 15 & 70 & 160 & 246 & 298 & 303\\
	\hline
	$b_{6,k}$ & 1 & 15 & 65 & 147 & 228 & 284 & 302 \\
	\hline
	$a_{6,k} \, \openbox \, b_{6,k}$ & $=$ & $=$ & $>$ & $>$ & $>$ & $>$ & $>$
\end{tabular}
\end{flushleft}
%\end{table}
%
%\begin{table}[h]
\begin{flushright}
\begin{tabular}{c||c|c|c|c|c|c|c|c|c||c}
	$k$ &7 & 8 &9 &10 &11 &12 &13 &14 &15 & \textbf{Total}\\
	\hline 
	$a_{6,k}$ &268 & 208 & 145 & 89 & 49 & 24 & 11 & 4 & 1 & \textbf{1892}\\
	\hline
	$b_{6,k}$ &277 & 223 & 157 & 98 & 53 & 26 & 11 & 4 & 1 & \textbf{1892}\\
	\hline
	$a_{6,k} \, \openbox \, b_{6,k}$ & $<$ & $<$ & $<$ & $<$ & $<$ & $<$ & $=$ & $=$ & $=$
\end{tabular}
\end{flushright}
%}
%\begin{table}[h]
\caption{$|\{\pi \in \mathfrak{D}^1_{12} (2143) | (2\text{-}31)\pi = k\}|$ versus $|\{\pi \in \mathfrak{D}^1_{12} (3421) | (13\text{-}2)\pi = k\}|$}
\label{table:cum_dist2}
\end{table}


%\vspace{0.1in}
\begin{table}[!h]
%{\small
\begin{flushleft}
\begin{tabular}{c||c|c|c|c|c|c|c|c|c}
	$k$  &0 & 1 & 2 & 3 & 4 & 5 & 6 & 7 & 8\\
	\hline 
	$a_{7,k}$ & 1 & 21 & 140 & 455 & 945 & 1497 & 1956 & 2215 & 2226\\
	\hline
	$b_{7,k}$ & 1 & 21 & 125 & 388 & 804 & 1294 & 1760 & 2089 & 2211 \\
	\hline
	$a_{7,k} \, \openbox \, b_{7,k}$ & $=$ & $=$ & $>$ & $>$ & $>$ & $>$ & $>$ & $>$ & $>$ 
\end{tabular}
\end{flushleft}
%\end{table}
%
%\begin{table}[h]
\begin{flushright}
\begin{tabular}{c||c|c|c|c|c|c|c|c|c|c|c|c|c||c}
	$k$  &9 &10 &11 &12 &13 &14 &15 &16 &17 &18 &19 &20 &21 & \textbf{Total}\\
	\hline 
	$a_{7,k}$ & 2032 & 1700 & 1317 & 948 & 641 & 410 & 249 & 140 & 72 & 32 & 13 & 4 & 1 & \textbf{17015}\\
	\hline
	$b_{7,k}$ & 2111 & 1840 & 1472 & 1091 & 750 & 482 & 288 & 158 & 78 & 34 & 13 & 4 & 1 & \textbf{17015}\\
	\hline
	$a_{7,k} \, \openbox \, b_{7,k}$ & $<$ & $<$ & $<$ & $<$ & $<$ & $<$ & $<$ & $<$ & $<$ & $<$ & $=$ & $=$ & $=$
\end{tabular}
\end{flushright}
%}
%\begin{table}[h]
\caption{$|\{\pi \in \mathfrak{D}^1_{14} (2143) | (2\text{-}31)\pi = k\}|$ versus $|\{\pi \in \mathfrak{D}^1_{14} (3421) | (13\text{-}2)\pi = k\}|$}
\label{table:cum_dist3}
\end{table}

%\newpage
\subsection{Conclusion}

In this paper, we have enumerated Dumont-4 permutations avoiding certain patterns, as well as Dumont-4 permutations containing a single occurrence of certain patterns. 
The Catalan numbers and powers of $2$ that occur in the enumeration of permutations avoiding three-letter patterns were also encountered in our results. However, for some patterns we could only find a generating function in the form of a continued fraction, as in \cite{RWZ}.

We also gave an intriguing conjecture regarding the Wilf-equivalence of a pair of patterns on Dumont permutations of the first kind. This conjecture has proven quite unyielding, and despite being presented twice at the open problem sessions at Permutation Patterns 2016 and 2018, progress on it is yet to be made. 

It would also be interesting to see if other Wilf-equivalences exist on Dumont permutations of the first kind avoiding a single four-letter pattern. The same question may be posed for the other kinds of Dumont permutations as well.

% End Section 5 - Conjecture and conclusion

\subsection*{Acknowledgments}

The authors would like to thank Michael Albert and Lindsey-Kay Lauderdale for computing several terms of sequences $|\mathfrak{D}^1_{2n} (2143)|$ and $|\mathfrak{D}^4_{2n} (312)|$, as well as the anonymous referee for a thorough reading of the paper and many helpful comments and suggestions.

%%%%%%%%%%%%%%%%%%%%

{\small

\begin{thebibliography}{99}

\bibitem{Albert_comm} M. Albert, personal communication, 2016.

%\bibitem{Albert} M. Albert, M. Atkinson, Simple permutations and pattern restricted permutations, \emph{Discrete Math.} \textbf{300} (2005), no. 1-3, 1-15.

\bibitem{Babson} E. Babson, E. Steingr\'imsson, Generalized permutation patterns and a classification of the Mahonian statistics, \emph{S\'em. Lothar. Combin.}, \textbf{44} (2000), Research Article B44b, 18 pp.

%\bibitem{Backelin} J. Backelin, J. West, G. Xin, Wilf equivalence for singleton classes. \emph{Proceedings of the 13th Conference on Formal Power Series and Algebraic Combinatorics}, Tempe, AZ, 2001.

\bibitem{Bevan} D. Bevan, \emph{Permutation patterns: basic definitions and notation}, \href{https://arxiv.org/abs/1506.06673}{arxiv:1506.06673}.

\bibitem{Bona_Burstein} M. B\'ona, A. Burstein, Permutations with exactly one copy of a decreasing pattern of length $k$, \emph{Ann. Combin.}, to appear, preprint at \href{https://arxiv.org/abs/2101.00332}{arxiv:2101.00332}.

\bibitem{Bona} M. B\'ona, \emph{Combinatorics of Permutations}. Chapman \& Hall/CRC, Boca Raton, FL, 2004.

%\bibitem{Bona_132} M. B\'ona, Permutations with one or two 132-subsequences, \emph{Discrete Math.} \textbf{181} (1998), 267-274.

%\bibitem{Bona_length4} M. B\'ona, Permutations avoiding certain patterns; The case of length 4 and generalizations.  \emph{Discrete Math.} \textbf {175} (1997), no. 1-3, 55-67.

\bibitem{Burstein_Short} A. Burstein, A Short Proof for the Number of Permutations Containing Pattern 321 Exactly Once, \emph{Electron. J. Combin.} \textbf{18} (2011).

%\bibitem {Burstein_Statistics} A. Burstein, On joint distribution of adjacencies, descents and some Mahonian statistics, \emph{Proceedings FPSAC 2010}, August 2010, San Francisco, CA, pp. 601-612.

%\bibitem {Bursten_Mahonian} A. Burstein, On the distribution of some Euler-Mahonian statistics, \emph{J. Combin.}, \textbf {6} (2015), 273-284.

\bibitem{Burstein_Restricted} A. Burstein, Restricted Dumont permutations, \emph{Ann. Combin.} \textbf{9} (2005), 269-280.

\bibitem{Burstein_Dyck} A. Burstein, S. Elizalde, T. Mansour, Restricted Dumont permutations, Dyck paths and noncrossing partitions, \emph{Discrete Math.} \textbf{306} (2006), no. 22, 2851-2869. Extended abstract in \emph{Proceedings of FPSAC 2006}, June 2006, San Diego, CA.

\bibitem{Burstein_Jones} A. Burstein, O. Jones, Restricted Dumont Permutations, talk, Special Session on Enumerative Combinatorics, AMS Central Section Meeting, University of St. Thomas, Minneapolis, MN, October 2016.

\bibitem{Burstein_Ofodile} A. Burstein, C. Ofodile, Dumont permutations containing one occurrence of certain three and four letter patterns, talk, Permutation Patterns 2011, California Polytechnic University, San Luis Obispo, CA, June 2011.

\bibitem {Burstein_New} A. Burstein, M. Josuat-Verg\`es, W. Stromquist, New Dumont permutations, \emph{Pure Math. Appl. (Pu.M.A.)} \textbf{21} (2010), no. 2, 177-206.

\bibitem{Burstein_Third} A. Burstein, W. Stromquist, Dumont permutations of the third kind, extended abstract, in \emph{Proceedings of FPSAC 2007}, July 2007, Tianjin, China.

\bibitem {Dumont} D. Dumont, Interpr\'etations combinatoires des nombres de Genocchi, \emph{Duke J. Math.} \textbf {41} (1974), 305-318.

\bibitem {Dumont_Foata} D. Dumont, D. Foata, Une propri\'et\'e de sym\'etrie des nombres de Genocchi, \emph{Bul. Soc. Math. France} \textbf{104} (1976), no. 4, 433-451.

\bibitem{Foata} Dominique Foata, On the Netto inversion number of a sequence, \emph{Proc. Amer. Math. Soc.} \textbf {19} (1968), 236-240.

%\bibitem {Gessel} I. Gessel, Symmetric functions and $p$-recursiveness, \emph{J. Combin. Th., Ser. A} \textbf {53} (1990), 257-285.

\bibitem {Jones} O. Jones, \emph{Enumerations of Dumont permutations avoiding certain four-letter patterns}, Ph.D. thesis, Howard University, 2019.

\bibitem {Kitaev_Remmel_1} S. Kitaev, J. Remmel, Classifying descents according to equivalence mod $k$, \emph{Electron. J. Combin.} \textbf{13(1)} (2006), \#R64.

\bibitem {Kitaev_Remmel_2} S. Kitaev, J. Remmel, Classifying descents according to parity, \emph{Ann. Combin.} \textbf{11} (2007), 173-193.

\bibitem{Krattenthaler} C. Krattenthaler, Permutations with restricted patterns and Dyck paths, \emph{Adv. Appl. Math.} \textbf{27} (2001), 510-530.

\bibitem {Knuth} D. E. Knuth, \emph{The Art of Computer Programming}, vol. 3. Addison-Wesley Publishing, MA 1973.

%\bibitem {MacMahon} P. MacMahon, \emph{Combinatory Analysis}, 2 volumes, Cambridge University Press, London, 1915-1916. Reprinted by Chelsea, New York, 1960.

\bibitem {Mansour} T. Mansour, Restricted 132-Dumont permutations, \emph{Australasian J. Combin.} \textbf{29} (2004), 103-117.

%\bibitem {MV_Restricted} T. Mansour,  A. Vainshtein, Restricted 132-avoiding permutations, \emph{Adv. Appl. Math.} \textbf{26} (2001), 258-269.

%\bibitem {MV_Chebyshev} T. Mansour, A. Vainshtein, Restricted Permutations and Chebyshev Polynomials, \emph{S\'em. Lothar. Combin.} \textbf{47} (2002), Article B47c.

\bibitem {Noonan} J. Noonan, The number of permutations containing exactly one increasing subsequence of length three, \emph{Discrete Math.} \textbf{152} (1996), no. 1-3, 307-313.

\bibitem{OEIS} OEIS Foundation Inc. (2021), \emph{The On-Line Encyclopedia of Integer Sequences}, \href{http://oeis.org}{http://oeis.org}.

\bibitem {Ofodile} C. Ofodile, \emph{The enumeration of Dumont permutations with few occurrences of three- and four-letter patterns}, Ph.D. thesis, Howard University, 2011.

\bibitem {Noonan_Zeilberger} J. Noonan, D. Zeilberger, The enumeration of permutations with a prescribed number of ``forbidden" patterns, \emph{Adv. Appl. Math.} \textbf{17} (1996), no. 4, 381-407.

\bibitem {Simion} R. Simion, F. W. Schmidt, Restricted permutations,  \emph{Europ. J. Combin.} \textbf {6} (1985), 383-406.

%\bibitem {Sloane} N.J.A. Sloane, S. Plouffe, \emph{The Encyclopedia of Integer Sequences}, Academic Press, New York (1995). Online at \href{http://oeis.org}{http://oeis.org}.

%\bibitem {Stankova} Z. Stankova, Forbidden subsequences. \emph{Discrete Math.} \textbf {132} (1994), no. 1-3, 291-316.

\bibitem{RWZ} A. Robertson, H.S. Wilf, D. Zeilberger, Permutation patterns and continued fractions, \emph{Electron. J. Combin.} \textbf{6} (1999), \#R38, 6 pp.

\bibitem{Vatter} V. Vatter, ``Permutation Classes'', in \emph{Handbook of Enumerative Combinatorics}, M. B\'ona, ed., Chapman and Hall/CRC, 2015. Available online at \href{https://arxiv.org/abs/1409.5159}{arxiv:1409.5159}.

%\bibitem {Wilf} H. Wilf, \emph{generatingfunctionology}, Academic Press, Boston (1994).

\bibitem {Zeilberger_lovely} D. Zeilberger, Alexander Burstein's lovely combinatorial proof of John Noonan's beautiful theorem that the number of $n$-permutations that contain the pattern $321$ exactly once equals $(3/n)(2n)!/((n-3)!(n+3)!)$, \emph{Pure Math. Appl.} \textbf{22} (2011), no. 2, 297-298.

\end{thebibliography}

}

\end{document}
